\documentclass[a4paper,11pt]{article}
% ---------- Кириллица: работает и в pdfLaTeX, и в XeLaTeX/LuaLaTeX ----------
\usepackage{iftex}
\ifPDFTeX
  \usepackage[T2A]{fontenc}
  \usepackage[utf8]{inputenc}
  \usepackage[russian]{babel}
\else
  \usepackage{fontspec}
  \IfFontExistsTF{CMU Serif}{\setmainfont{CMU Serif}}{%
    \IfFontExistsTF{Times New Roman}{\setmainfont{Times New Roman}}{\setmainfont{FreeSerif}}}
  \usepackage{polyglossia}
  \setdefaultlanguage{russian}
\fi
\usepackage{amsmath,amssymb,mathtools}
\usepackage[a4paper,margin=2cm]{geometry}
\usepackage[dvipsnames]{xcolor}
\usepackage{tikz}
\usetikzlibrary{arrows.meta,calc,patterns,decorations.pathreplacing,babel}
\usepackage[most]{tcolorbox}
\usepackage{enumitem}
\usepackage{titlesec}
\usepackage{array}
\usepackage[unicode,hidelinks]{hyperref}

\emergencystretch=3em
\tolerance=1500
\setlength{\parindent}{0pt}
\setlength{\parskip}{4pt}
\setlist{nosep,leftmargin=*,topsep=2pt}

% ---------- оформление ----------
\definecolor{secred}{RGB}{190,40,45}
\definecolor{subteal}{RGB}{0,110,120}
\definecolor{leadcol}{RGB}{0,120,110}
\titleformat{\section}{\Large\bfseries\color{secred}}{\thesection.}{0.5em}{}
\titleformat{\subsection}{\large\bfseries\color{subteal}}{\thesubsection.}{0.5em}{}
\titlespacing*{\section}{0pt}{14pt}{6pt}
\titlespacing*{\subsection}{0pt}{10pt}{4pt}

\newtcolorbox{opr}{breakable,enhanced,colback=blue!4,colframe=blue!55!black,
  boxrule=0pt,leftrule=2.2pt,arc=0pt,left=6pt,right=4pt,top=3pt,bottom=3pt,
  before skip=6pt,after skip=6pt}
\newtcolorbox{teor}[1]{breakable,enhanced,colback=orange!7,colframe=orange!70!black,
  boxrule=0pt,leftrule=2.2pt,arc=0pt,left=6pt,right=4pt,top=3pt,bottom=3pt,
  before skip=6pt,after skip=6pt,before upper={\textbf{Теорема #1.}\ }}
\newtcolorbox{hwbox}{breakable,enhanced,colback=gray!6,colframe=gray!60,
  boxrule=0.4pt,arc=1pt,left=6pt,right=6pt,top=3pt,bottom=3pt,fontupper=\small}

\newcommand{\lead}[1]{\textbf{\color{leadcol}#1}\ }
\newcommand{\q}{\textbf{\color{red!70!black}?}\ }
\newcommand{\imp}{\textbf{\color{red!70!black}!}\ }
\newcommand{\zam}{\textbf{Замечание.}\ }
\newcommand{\prim}{\textbf{Пример.}\ }
\newcommand{\ispr}[1]{\footnote{\textit{Исправлено.} #1}}
\renewcommand{\thefootnote}{\textcolor{red!75!black}{\arabic{footnote}}}
\newcommand{\hwlist}[1]{{\small\color{gray}[#1]}}
\newcommand{\qedr}{\hfill$\blacksquare$}
\DeclareRobustCommand{\cn}[1]{\tikz[baseline=(c.base)]\node[draw,circle,inner sep=0.7pt,minimum size=1.15em,font=\footnotesize\bfseries] (c) {#1};}

\newcommand{\rtg}{\operatorname{tg}}
\newcommand{\rarctg}{\operatorname{arctg}}
\newcommand{\Rez}{\operatorname{Re}}
\newcommand{\Imz}{\operatorname{Im}}
\newcommand{\rang}{\operatorname{rang}}
\newcommand{\diag}{\operatorname{diag}}
\newcommand{\gauss}[2][]{\;\xrightarrow[\scriptstyle #1]{\scriptstyle #2}\;}


\begin{document}
\begin{center}
{\LARGE\bfseries\color{secred}Аналитическая геометрия}\\[3pt]
{\large Конспект лекций}
\end{center}
\vspace{-2mm}

\tableofcontents
\newpage


% =====================================================================
\section{Матрицы}
% =====================================================================
Матрицы возникают как средство решения линейных систем.

\prim Предприятие использует $m$ ресурсов в объёмах $b_1,\dots,b_m$, выпуская $n$ продуктов
в количествах $x_1,\dots,x_n$. Пусть $a_{ij}$~--- расход ресурса $i$ на единицу продукта $j$
(например, $a_{m2}$~--- использование ресурса $m$ на выпуск продукта \textnumero\,2, а $x_2$~--- объём
этого продукта). Тогда $x_j$ удовлетворяют системе
\[
\begin{cases}
a_{11}x_1+a_{12}x_2+\dots+a_{1n}x_n=b_1,\\
\quad\dots\dots\dots\dots\dots\dots\dots\dots\dots\\
a_{m1}x_1+a_{m2}x_2+\dots+a_{mn}x_n=b_m,
\end{cases}
\qquad x_1,\dots,x_n\ \text{~--- ?}
\]

\lead{Идея:} решать систему, используя таблицу её коэффициентов.

\begin{opr}
\textbf{Определение.} \emph{Матрица} $A$ размерности $m\times n$~--- это таблица из $m$ строк
и $n$ столбцов с элементами $a_{ij}$ ($i$~--- номер строки, $j$~--- номер столбца).
\end{opr}

\prim $A=\begin{pmatrix}2&3&0\\-1&4&5\end{pmatrix}$~--- матрица $2\times3$, \ $a_{13}=0$, \ $a_{22}=4$.

\subsection{Частные случаи матриц}
\begin{enumerate}[label=\arabic*)]
\item \textbf{Строка} (вектор-строка): $a=(1\ \ 2\ \ 3)$~--- матрица $1\times3$, в общем случае $1\times n$.
\item \textbf{Столбец} (вектор-столбец): $a=\begin{pmatrix}-1\\2\\0\end{pmatrix}$~--- матрица $3\times1$.

\imp Далее \emph{вектор}~--- всегда столбец, т.\,е. упорядоченный набор чисел (матрица $m\times1$).

\imp Матрицы обозначаем большими буквами, векторы~--- маленькими.

\item \textbf{Треугольная матрица}: все элементы выше (ниже) главной диагонали равны нулю:
$a_{ij}=0$ при $i<j$ (при $i>j$).
\[
\underbrace{\begin{pmatrix}\ast&&0\\ \ast&\ast&\\ \ast&\ast&\ast\end{pmatrix}}_{\text{нижняя}}
\qquad
\underbrace{\begin{pmatrix}\ast&\ast&\ast\\ &\ast&\ast\\ 0&&\ast\end{pmatrix}}_{\text{верхняя}}
\qquad
\prim \begin{pmatrix}1&4&5\\0&2&6\\0&0&3\end{pmatrix}
\]
Элементы $a_{ii}$ образуют (главную) диагональ; в последнем примере под диагональю нули~---
это верхняя треугольная матрица.

\item \textbf{Диагональная матрица}: $a_{ij}=0$ при $i\ne j$. \
\prim $\begin{pmatrix}0&0&0&0\\0&1&0&0\\0&0&2&0\end{pmatrix}$~--- элементы с $i=j$ тоже могут быть равны нулю.

\imp \emph{Единичная матрица} $E=\begin{pmatrix}1&&0\\&\ddots&\\0&&1\end{pmatrix}$ ($n\times n$)~---
и диагональная, и треугольная.

\imp \emph{Нулевая матрица} $O_{mn}$ (все элементы~--- нули)~--- и треугольная, и диагональная.
\end{enumerate}

\subsection{Операции над матрицами}
\begin{enumerate}[label=\arabic*)]
\item \textbf{Транспонирование.} Каждая строка \textnumero\,$i$ матрицы $A$ становится столбцом
\textnumero\,$i$ матрицы $A^T$ (строки и столбцы <<меняются местами>>). При этом $A$
<<поворачивается>> относительно диагонали:
\[
\begin{pmatrix}1&2&3&4\\5&6&7&8\end{pmatrix}^{T}=\begin{pmatrix}1&5\\2&6\\3&7\\4&8\end{pmatrix}
\quad\text{(диагональ $1,6$ остаётся на месте).}
\]

\item \textbf{Линейные операции:}
 \begin{enumerate}[label=\alph*)]
 \item умножение на число: $\alpha A=(\alpha a_{ij})$~--- каждый элемент умножается на число;
 \item сложение: $A+B=C$, \ $c_{ij}=a_{ij}+b_{ij}$~--- складываются соответствующие элементы,
 поэтому размерности $A$ и $B$ должны совпадать.
 \end{enumerate}

\item \textbf{Скалярное произведение векторов} (алгебраическое обозначение):
\[
a^Tb=(a_1\ \dots\ a_n)\begin{pmatrix}b_1\\ \vdots\\ b_n\end{pmatrix}=\sum_{i=1}^{n}a_ib_i
\ \text{~--- сумма произведений соответствующих координат.}
\]
\zam Геометрически $a^Tb=0\iff a\perp b$; \ $a^2=a^Ta$~--- скалярный квадрат; \
$|a|=\sqrt{a^2}$~--- евклидова длина вектора $a$.

\item \textbf{Произведение матриц} $C=AB$: если $A$~--- $n\times k$, $B$~--- $k\times m$, то $C$~--- $n\times m$,
\[
c_{ij}=\text{(строка \textnumero\,$i$ матрицы $A$)}\cdot\text{(столбец \textnumero\,$j$ матрицы $B$)}
\ \text{~--- скалярное произведение.}
\]
$\Rightarrow$ Столбцов в левой матрице должно быть столько же, сколько строк в правой.

\textbf{Примеры.}
\begin{enumerate}[label=\arabic*)]
\item $\begin{pmatrix}1\\2\\3\end{pmatrix}(4\ \ 5\ \ 6)$: \ $3\times1$ на $1\times3$~--- можно, получится $3\times3$:
$\begin{pmatrix}1\cdot4&1\cdot5&1\cdot6\\2\cdot4&2\cdot5&2\cdot6\\3\cdot4&3\cdot5&3\cdot6\end{pmatrix}
=\begin{pmatrix}4&5&6\\8&10&12\\12&15&18\end{pmatrix}$.
\item $\begin{pmatrix}1\\2\end{pmatrix}\begin{pmatrix}3\\4\end{pmatrix}$: \ $2\times1$ на $2\times1$, \ $1\ne2$~---
перемножать нельзя.
\item $\begin{pmatrix}1&2\\3&4\end{pmatrix}\begin{pmatrix}5&7&9\\6&8&10\end{pmatrix}$: \ $2\times2$ на $2\times3$~---
можно, <<строка на столбец>>:
\[
=\begin{pmatrix}1\cdot5+2\cdot6&1\cdot7+2\cdot8&1\cdot9+2\cdot10\\3\cdot5+4\cdot6&3\cdot7+4\cdot8&3\cdot9+4\cdot10\end{pmatrix}
=\begin{pmatrix}17&23&29\\39&53&67\end{pmatrix}\quad(2\times3).
\]
\end{enumerate}
\imp Произведение $BA$ здесь найти нельзя ($2\times3$ на $2\times2$, \ $3\ne2$). В общем случае
$AB\ne BA$. Если $AB=BA$, то говорят, что $A$ и $B$ \emph{коммутируют}.
\end{enumerate}

\subsection{Свойства операций}
В основном они аналогичны свойствам операций с числами и непосредственно следуют из определений.

\textbf{Сложение:} $A+B=B+A$ (коммутативность); \ $A+(-A)=O$ (противоположный элемент); \
$A+O=A$ (нейтральный элемент~--- нулевая матрица); \ $A+(B+C)=(A+B)+C$ (ассоциативность); \
$(A+B)^T=A^T+B^T$.

\textbf{Умножение на число:} $\alpha(A+B)=\alpha A+\alpha B$, \ $(\alpha+\beta)A=\alpha A+\beta A$
(дистрибутивность); \ $(\alpha\beta)A=\alpha(\beta A)$ (ассоциативность).

\textbf{Умножение матриц:} $C(A+B)=CA+CB$, \ $(A+B)C=AC+BC$ (дистрибутивность)\ispr{В рукописи
записано $c(A+B)=(A+B)c=cA+cB$. Для матрицы $C$ порядок множителей важен ($AB\ne BA$), поэтому
это два разных равенства.}; \
$A(BC)=(AB)C$ (ассоциативность); \ $\alpha(AB)=(\alpha A)B=A(\alpha B)$; \
$AE=EA=A$ (нейтральность $E$); \ $(AB)^T=B^TA^T$.

% =====================================================================
\section{Определитель матрицы (детерминант)}
% =====================================================================
Определитель вычисляется для квадратной матрицы.
\begin{enumerate}[label=\arabic*)]
\item $A=(a_{11})$~--- матрица $1\times1$: \ $\det A=a_{11}$.
\item $A$~--- матрица $2\times2$:
$\det A=\begin{vmatrix}a_{11}&a_{12}\\a_{21}&a_{22}\end{vmatrix}=a_{11}a_{22}-a_{12}a_{21}$
(произведение диагональных элементов минус произведение внедиагональных).
\end{enumerate}

\begin{opr}
\textbf{Определение.} Определитель квадратной $n\times n$-матрицы~--- это скалярное произведение
любой строки (столбца) на вектор алгебраических дополнений элементов этой строки (столбца):
\[
\det A=\sum_{j=1}^{n}a_{ij}A_{ij}\ \ \text{~--- разложение по строке $i$},\qquad
\det A=\sum_{i=1}^{n}a_{ij}A_{ij}\ \ \text{~--- разложение по столбцу $j$}.
\]
Здесь $A_{ij}=(-1)^{i+j}M_{ij}$~--- \emph{алгебраическое дополнение} элемента $a_{ij}$;
$M_{ij}$~--- \emph{минор} элемента $a_{ij}$: определитель матрицы, полученной из $A$ вычёркиванием
строки $i$ и столбца $j$.
\end{opr}

\prim Разложение по 3-му столбцу\ispr{В рукописи в ответе стоит 16; проверено вычислением: определитель равен~7.}:
\[
\begin{vmatrix}5&-1&4\\3&0&1\\2&3&-4\end{vmatrix}
=4\cdot(-1)^{1+3}\underbrace{\begin{vmatrix}3&0\\2&3\end{vmatrix}}_{9}
+1\cdot(-1)^{2+3}\underbrace{\begin{vmatrix}5&-1\\2&3\end{vmatrix}}_{17}
+(-4)\cdot(-1)^{3+3}\underbrace{\begin{vmatrix}5&-1\\3&0\end{vmatrix}}_{3}
=36-17-12=7.
\]

Для определителя 3-го порядка существуют также:
\begin{enumerate}[label=\arabic*)]
\item \textbf{правило треугольников}:
\[
\det A=a_{11}a_{22}a_{33}+a_{12}a_{23}a_{31}+a_{21}a_{32}a_{13}-a_{13}a_{22}a_{31}-a_{12}a_{21}a_{33}-a_{23}a_{32}a_{11};
\]
\begin{center}
\begin{tikzpicture}[scale=0.55]
 \foreach \i in {1,2,3}{\foreach \j in {1,2,3}{\fill (\j,-\i) circle (2.5pt);}}
 \draw[thick,blue] (1,-1)--(3,-3);
 \draw[thick,blue] (2,-1)--(3,-2)--(1,-3)--cycle;
 \draw[thick,blue] (1,-2)--(2,-3)--(3,-1)--cycle;
 \node at (2,-4) {<<$+$>>};
 \begin{scope}[xshift=5cm]
 \foreach \i in {1,2,3}{\foreach \j in {1,2,3}{\fill (\j,-\i) circle (2.5pt);}}
 \draw[thick,red] (3,-1)--(1,-3);
 \draw[thick,red] (2,-1)--(1,-2)--(3,-3)--cycle;
 \draw[thick,red] (2,-3)--(3,-2)--(1,-1)--cycle;
 \node at (2,-4) {<<$-$>>};
 \end{scope}
\end{tikzpicture}
\end{center}
\item \textbf{правило Саррюса}: справа приписываем 1-й и 2-й столбцы; три <<диагонали>> вниз-вправо
берём со знаком <<$+$>>, три вниз-влево~--- со знаком <<$-$>>:
\[
\det A=a_{11}a_{22}a_{33}+a_{12}a_{23}a_{31}+a_{13}a_{21}a_{32}-a_{13}a_{22}a_{31}-a_{11}a_{23}a_{32}-a_{12}a_{21}a_{33}.
\]
\begin{center}
\begin{tikzpicture}[scale=0.85]
 \foreach \i in {1,2,3}{\foreach \j in {1,...,5}{
   \pgfmathtruncatemacro{\jj}{\j>3 ? \j-3 : \j}
   \node (m\i\j) at (\j*1.1,-\i*0.7) {\small$a_{\i\jj}$};}}
 \draw[gray] (3.65,-0.35)--(3.65,-2.45);
 \foreach \s in {1,2,3}{
   \pgfmathtruncatemacro{\b}{\s+2}
   \draw[blue,thick,opacity=0.6] (m1\s.north west)--(m3\b.south east);}
 \foreach \s in {3,4,5}{
   \pgfmathtruncatemacro{\b}{\s-2}
   \draw[red,thick,opacity=0.6] (m1\s.north east)--(m3\b.south west);}
\end{tikzpicture}
\end{center}
\end{enumerate}

\subsection{Свойства определителей}
\begin{enumerate}[label=\arabic*)]
\item $|A^T|=|A|$~--- при транспонировании определитель сохраняется.
\item $|AB|=|A|\cdot|B|$.
\item Если все элементы строки (столбца) умножить на число, то определитель умножится на это число.
\emph{Следствия:} а) если есть нулевая строка (столбец), то $|A|=0$; \ б) $|\lambda A|=\lambda^n|A|$
($A$~--- $n\times n$).
\item При перестановке двух строк (столбцов) определитель меняет знак.
\item Для треугольной матрицы $\det A=\prod_{i=1}^{n}a_{ii}$~--- произведение диагональных элементов.
\item Определитель не изменится при \emph{элементарном преобразовании}: к элементам строки (столбца)
прибавить элементы другой строки (столбца), умноженные на число.
\emph{Следствие:} если две строки (столбца) равны или пропорциональны, то определитель равен нулю.
\end{enumerate}

\subsection{Обоснование свойств определителей}
Если разложить определитель до конца, то окажется
\[
|A|=\sum_{\sigma}(-1)^{|\sigma|}a_{1\sigma(1)}a_{2\sigma(2)}\dots a_{n\sigma(n)}
=\sum_{\sigma}(-1)^{|\sigma|}\prod_{i=1}^{n}a_{i\sigma(i)}\tag{1}
\]
~--- сумма по всем перестановкам $\sigma$ множества $\{1,2,\dots,n\}$. Всего $n!$ вариантов
$\Rightarrow$ $n!$ слагаемых.

$|\sigma|$~--- \emph{число инверсий} перестановки: количество случаев, когда $i<j$, но $\sigma(i)>\sigma(j)$.

\prim $\{1,2,3,4\}\to\{4,2,1,3\}$: инверсии $(4,2),\ (4,1),\ (4,3),\ (2,1)$~--- их 4, перестановка
<<чётная>>.\ispr{В рукописи выписаны только две инверсии: $(4,2)$ и $(2,1)$; вывод о чётности от этого
не меняется.}

\begin{enumerate}[label=\arabic*.]
\item $A$~--- $2\times2$: по (1) $|A|=(-1)^0a_{11}a_{22}+(-1)^1a_{12}a_{21}$ (тождественная
перестановка и перестановка $\{1,2\}\to\{2,1\}$).
\item Определитель порядка 3~--- см. в методичке.
\end{enumerate}

\textbf{Установим свойства.}
\begin{enumerate}[label=\arabic*)]
\item Если строку (столбец) умножить на $\lambda$, то $\lambda$ выносится за знак определителя.
Действительно, по (1)
$|B|=\sum_\sigma(-1)^{|\sigma|}a_{1\sigma(1)}\cdots\lambda a_{i\sigma(i)}\cdots a_{n\sigma(n)}=\lambda|A|$.

Если в определителе нулевая строка (столбец), то он равен 0: можно считать, что вся строка
имеет множитель $\lambda=0$.

\item $\det A^T=\det A$. Оба определителя вычисляются по (1) и имеют одинаковые слагаемые:
$\det A^T=\sum_\sigma(-1)^{|\sigma|}\prod_{i}a_{\sigma(i)i}$~--- те же слагаемые, что и в (1).
Самостоятельно проверить для $2\times2$, $3\times3$.

\item При перестановке двух строк (столбцов) определитель меняет знак: все слагаемые те же, но
чётность каждой перестановки меняется на 1 $\Rightarrow$ все $(-1)^{|\sigma|}$ меняют знак
$\Rightarrow$ получаем $-\det A$.

\textbf{Следствие!} Если в матрице две строки (столбца) пропорциональны или равны, то $\det A=0$.

Без ограничения общности пропорциональны строки 1 и 2:
\[
A=\begin{pmatrix}a_{11}&\dots&a_{1n}\\ \lambda a_{11}&\dots&\lambda a_{1n}\\ a_{31}&\dots&a_{3n}\\
\vdots&&\vdots\\ a_{n1}&\dots&a_{nn}\end{pmatrix}.
\]
Множитель $\lambda$ выносим за знак (свойство 1): $\det A=\lambda\det A'$. Переставим строки 1 и 2
$A'$~--- знак должен поменяться, но матрица не изменилась: $\det A'=-\det A'$ $\Rightarrow$ $\det A'=0$.

\item[4)!] Если к строке (столбцу) прибавить другую строку (столбец), умноженную на число,
определитель не изменится.

Без ограничения общности строку 1 умножим на $\lambda$ и прибавим к строке 2~--- получим матрицу $B$.
Разложим $\det B$ по строке 2:
\[
\det B=\sum_{j}(-1)^{2+j}(a_{2j}+\lambda a_{1j})M_{2j}
=\underbrace{\sum_{j}(-1)^{2+j}a_{2j}M_{2j}}_{=\det A}
+\lambda\underbrace{\sum_{j}(-1)^{2+j}a_{1j}M_{2j}}_{=0}=\det A,
\]
второе слагаемое равно нулю, так как это определитель матрицы, в которой одинаковы строки 1 и 2.

В частности, видим: \emph{скалярное произведение строки на вектор алгебраических дополнений любой
другой строки равно нулю} (то же для столбцов).
\end{enumerate}

% =====================================================================
\section{Обратная матрица}
% =====================================================================
\lead{Идея.} Линейное уравнение $ax=b$ имеет единственное решение, если $a\ne0$: \ $x=\frac1a\cdot b$,
где $\frac1a$~--- обратный элемент для $a$ при умножении.

\q Как обобщить это на матричное уравнение
\[
AX=B\qquad(A,\ B\ \text{даны},\ X\ \text{найти})?\tag{1}
\]

\begin{opr}
\textbf{Определение.} $A^{-1}$~--- \emph{обратная матрица} для $A$, если
\[
A^{-1}A=AA^{-1}=E\tag{2}
\]
($E$~--- единичная матрица, нейтральный элемент при умножении матриц).
\end{opr}

\zam\imp В математике любое определение должно быть \emph{целесообразным} (предназначено для решения
какой-то задачи) и \emph{корректным}, т.\,е. применимым к каким-то объектам.

\q Целесообразно ли определение $A^{-1}$? Пусть $\exists A^{-1}$; умножим (1) на $A^{-1}$ слева:
$A^{-1}AX=A^{-1}B$, а по (2) $A^{-1}AX=EX=X$ $\Rightarrow$ (1) имеет единственное решение
\[
X=A^{-1}B.\tag{3}
\]

\q Корректно ли? Да: $A^{-1}$ существует для \emph{невырожденной} матрицы (определитель $\ne0$).
\begin{enumerate}[label=\arabic*)]
\item \textbf{Необходимость.} По свойству определителей
$\det(AA^{-1})=\det A\cdot\det A^{-1}=\det E=1$ $\Rightarrow$ для существования $A^{-1}$ необходимо
\[
\det A\ne0.\tag{4}
\]
\item \textbf{Достаточность.} $A^{-1}$ существует и вычисляется по формуле
\[
A^{-1}=\frac{1}{\det A}\,\tilde A^{T},\tag{5}
\]
где $\tilde A^T$~--- транспонированная матрица алгебраических дополнений элементов $A$
(\emph{присоединённая матрица}).

Действительно, $A^{-1}A=\frac{1}{|A|}\tilde A^TA$; строки левой матрицы $\tilde A^T$~--- алгебраические
дополнения элементов столбцов $A$. В силу замечания в конце \S\,2 при перемножении
$\tilde A^T\cdot A$ <<строка на столбец>> строка $i$ на столбец $j$ даёт $\det A$ при $j=i$ (это разложение
по столбцу $i$) или $0$ при $j\ne i$. В итоге
\[
A^{-1}A=\frac{1}{|A|}\begin{pmatrix}|A|&&0\\&\ddots&\\0&&|A|\end{pmatrix}=E.
\]
\end{enumerate}

\subsection{Метод Гаусса для обратной матрицы}
\textbf{Элементарные преобразования строк:} 1) умножение строки на число $\lambda\ne0$;
2) перестановка двух строк; 3) умножение строки на число и прибавление к другой строке.

Их можно реализовать с помощью левого матричного умножения $P\cdot A$:
\begin{itemize}
\item $P=\diag(1,\dots,\lambda,\dots,1)$ ($p_{ii}=\lambda$)~--- строка $i$ умножится на число $\lambda$;
\item $P=\begin{pmatrix}0&1&&\\1&0&&\\&&1&\\&&&\ddots\end{pmatrix}$~--- перестановка строк 1 и 2.
\q Какой вид имеет $P$ при перестановке произвольных строк?
\item к строке 1 прибавляется строка 2, умноженная на число $\lambda$: \
$P=\begin{pmatrix}1&\lambda&&\\0&1&&\\&&\ddots&\\&&&1\end{pmatrix}$.
\q Какой вид имеет $P$ для произвольных строк?
\end{itemize}
(Ответы на эти вопросы~--- в п.\,\ref{ssec:obrgauss}.)

\lead{Идея!} Пусть в результате элементарных преобразований строк из $(A\,|\,E)$ получилось $(E\,|\,D)$.
Это означает, что существуют матрицы $P_1,\dots,P_k$:
\[
P_kP_{k-1}\cdots P_1A=E,\qquad D=P_k\cdots P_1E=P_k\cdots P_1\ \Rightarrow\ DA=E\ \Rightarrow\ D=A^{-1}.
\]

% =====================================================================
\section{Решение СЛАУ методом Гаусса}
% =====================================================================
Рассмотрим систему линейных алгебраических уравнений (СЛАУ)
\[
a_{i1}x_1+a_{i2}x_2+\dots+a_{im}x_m=b_i,\quad i=1,\dots,n
\quad\iff\quad Ax=b,\tag{1}
\]
где $A=(a_{ij})$~--- основная матрица ($n\times m$), $x=(x_1,\dots,x_m)^T$~--- вектор неизвестных,
$b=(b_1,\dots,b_n)^T$~--- вектор свободных членов.

\imp\lead{Проблема.} Если $n=m$ (число уравнений равно числу неизвестных), то при $\det A\ne0$
система (1) имеет решение $x=A^{-1}b$ (прошлая лекция). Если же
\begin{enumerate}[label=\arabic*)]
\item $n>m$~--- система переопределена;
\item $n<m$~--- система недоопределена;
\item $n=m$, но $\det A=0$,
\end{enumerate}
то метод обратной матрицы <<не работает>> $\Rightarrow$ нужен новый метод решения СЛАУ.

\lead{Идея:} обобщим <<школьный метод>>: из одного уравнения выражаем переменную и подставляем
в другие, т.\,е. по сути во втором уравнении исключается одна из переменных $\Rightarrow$ обобщим
этот способ исключения переменных.

Средства реализации идеи~--- преобразования уравнений, при которых получается равносильная, но
более простая система (решения сохраняются). Составим для (1) \emph{расширенную матрицу
коэффициентов} $\bar A=(A\,|\,b)$ (к основной матрице справа добавлен столбец свободных членов). Тогда
\begin{enumerate}[label=\arabic*)]
\item умножить уравнение \textnumero\,$i$ на число $\lambda\ne0$ $\iff$ умножить на $\lambda$ строку
\textnumero\,$i$ матрицы $\bar A$;
\item переставить уравнения \textnumero\,$s$ и \textnumero\,$l$ $\iff$ переставить строки \textnumero\,$s$
и \textnumero\,$l$ в $\bar A$;
\item уравнение \textnumero\,$i$ умножить на $\lambda$ и сложить с уравнением \textnumero\,$s$ $\iff$
в $\bar A$ к строке \textnumero\,$s$ прибавить строку \textnumero\,$i$, умноженную на $\lambda$.
\end{enumerate}
Очевидно, это те же преобразования, что и в методе Гаусса для обратной матрицы, но там
преобразовывали $(A\,|\,E)$, а здесь~--- $(A\,|\,b)$.

\textbf{Цель преобразований:} исключить $x_1$ из уравнения 2, $x_1,x_2$~--- из уравнения 3 и т.\,д.:
в каждом уравнении \textnumero\,$i\geqslant2$ должны отсутствовать $x_1,\dots,x_{i-1}$.
В итоге получаем алгоритм метода Гаусса (пока не ясно, всегда ли он будет работать).

\subsection{Алгоритм метода Гаусса}
\textbf{\cn{1}\ Прямой ход метода.} Составив $\bar A$ для (1), с помощью элементарных
преобразований строк приводим её к \emph{ступенчатой форме} (верхней треугольной): под диагональю
нули, так как переменные исключены.

\textbf{\cn{2}\ Проверим, совместна ли система} (имеет ли она решения). Допустим,
в преобразованной матрице $\bar A'$
\begin{enumerate}[label=\arabic*)]
\item есть строка
$\left(\begin{array}{cccc|c}0&0&\dots&0&b_i'\end{array}\right)$, \ $b_i'\ne0$ (строка \textnumero\,$i$).
Ей соответствует преобразованное уравнение $0\cdot x_1+0\cdot x_2+\dots+0\cdot x_m=b_i'$~---
противоречие. В таком случае строка \textnumero\,$i$ \emph{противоречивая}.

\textbf{Вывод:} преобразованная равносильная система содержит противоречие $\Rightarrow$ система (1)
\emph{несовместна}.
\item есть строка $\left(\begin{array}{cccc|c}0&0&\dots&0&0\end{array}\right)$~--- строка
\textnumero\,$i$ \emph{нулевая} $\Rightarrow$ уравнение $0=0$ верно при любых $x$ $\Rightarrow$ его можно
исключить, т.\,е. такую строку из $\bar A'$ можно вычеркнуть (это уравнение~--- следствие других
уравнений).
\item Если в $\bar A'$ нет противоречивых строк, то система (1) \emph{совместна}.
\end{enumerate}

\textbf{\cn{3}\ Обратный ход метода Гаусса.} По преобразованной матрице $\bar A'$
составим равносильную систему (количество уравнений может стать $k\leqslant n$):
\[
\begin{cases}
a'_{11}x_1+a'_{12}x_2+\dots+a'_{1m}x_m=b'_1,\\
\phantom{a'_{11}x_1+{}}a'_{22}x_2+\dots+a'_{2m}x_m=b'_2,\\
\phantom{a'_{11}x_1+a'_{12}x_2+{}}\dots\dots\dots\dots\\
\phantom{a'_{11}x_1+a'_{12}x_2+{}}a'_{kk}x_k+\dots+a'_{km}x_m=b'_k.
\end{cases}\tag{2}
\]
Из последнего уравнения выразим $x_k$, подставим в предыдущее, найдём $x_{k-1}$ и т.\,д.

\textbf{Вывод~--- возможные варианты:}
\begin{enumerate}[label=\arabic*)]
\item Если в уравнении \textnumero\,$k$ системы (2) есть <<свободные переменные>>, то система (1)
имеет бесконечно много решений.

\prim $\begin{cases}x_1+x_2+x_3=1,\\ x_2-x_3=2\end{cases}$\ispr{В рукописи второе уравнение записано как
$x_2+x_3=2$, но дальше получено $x_2=2+x_3$ и $x_1=-1-2x_3$. Эти выражения верны для уравнения
$x_2-x_3=2$, поэтому исправлено условие.}
$\Rightarrow$ $x_2=2+x_3$, $x_3$~--- \emph{свободная} переменная, т.\,к. на неё никаких ограничений
нет; $x_1=1-x_2-x_3=-1-2x_3$:
\[
\begin{cases}x_1=-1-2x_3,\\x_2=2+x_3\end{cases}\quad\text{~--- бесконечно много решений.}
\]
Итак, если преобразованная система (2) недоопределена, т.\,е. $k<m$, то (1) имеет бесконечно много
решений.
\item Если $k=m$ в (2), то система имеет единственное решение (в этом случае изначально можно
было применить метод обратной матрицы).
\end{enumerate}

\subsection{Обоснование метода Гаусса}
\lead{Проблемы:} 1) неочевидно, что любую матрицу можно привести к ступенчатой форме;
2) описание признака совместности пока недостаточно математично.

\begin{teor}{1}
Любую матрицу с помощью элементарных преобразований строк можно привести к ступенчатому виду.
\end{teor}

\textbf{Схема доказательства.}
\begin{enumerate}[label=\arabic*)]
\item Рассмотрим столбец \textnumero\,1 матрицы $B=(b_{ij})$:
 \begin{enumerate}[label=\alph*)]
 \item все его элементы равны 0~--- значит, и ниже диагонали нули $\Rightarrow$ переходим к следующему
 столбцу;
 \item $\exists\,b_{i1}\ne0$~--- переставим строки $i$ и 1, получим $B'=(b'_{ij})$, $b'_{11}=b_{i1}\ne0$.
 Умножим строку 1 на $-\dfrac{b'_{i1}}{b'_{11}}$ и сложим со строкой $i$ ($i=2,3,\dots,n$) $\Rightarrow$
 в первом столбце ниже диагонали получим нули:
 $\begin{pmatrix}b'_{11}&\dots&\dots\\0&b'_{22}&\dots\\ \vdots&&\\0&\dots&\end{pmatrix}$.
 Далее переходим к следующему столбцу.
 \end{enumerate}
\item Аналогично рассмотрим столбцы $j=2,3,\dots,m$ преобразованной матрицы:
 \begin{enumerate}[label=\alph*)]
 \item все элементы столбца $j$ ниже диагонали равны 0 ($b_{ij}=0$, $i>j$) $\Rightarrow$ переходим
 к следующему столбцу;
 \item $\exists\,b_{ij}\ne0$, $i>j$~--- переставляем строки $i$ и $j$, получаем $B''=(b''_{ij})$,
 $b''_{jj}\ne0$ $\Rightarrow$ умножаем строку \textnumero\,$j$ на $-\dfrac{b''_{ij}}{b''_{jj}}$ и складываем
 с каждой строкой \textnumero\,$i>j$. В итоге под диагональю в столбце \textnumero\,$j$~--- нули.
 \end{enumerate}
\end{enumerate}
В итоге получим ступенчатую матрицу.\qedr

Введём понятия для проверки совместности.
\begin{opr}
\textbf{Определение.} Векторы $a_1,\dots,a_m$ \emph{линейно зависимы} (ЛЗ), если они могут
образовать нулевую линейную комбинацию с ненулевыми коэффициентами, т.\,е. существуют числа
$\alpha_1,\dots,\alpha_m$, не все равные нулю, такие, что
\[
\sum_{i=1}^{m}\alpha_ia_i=0\qquad\text{(линейная комбинация).}\tag{3}
\]
\end{opr}
Если в (3) есть $\alpha_i\ne0$, то один из векторов~--- линейная комбинация остальных: при $\alpha_1\ne0$
\ $a_1=-\sum\limits_{i=2}^{m}\dfrac{\alpha_i}{\alpha_1}a_i$.

\begin{opr}
\textbf{Определение!} Векторы $a_1,\dots,a_m$ \emph{линейно независимы} (ЛНЗ), если (3) верно лишь
при нулевых коэффициентах $\alpha_1=\dots=\alpha_m=0$, т.\,е. ни один из векторов нельзя выразить
линейной комбинацией других.
\end{opr}

\textbf{Применим эти определения к методу Гаусса.} Если в результате элементарных преобразований строк
получилась нулевая строка, то есть набор линейно зависимых строк $\iff$ одно из уравнений~---
следствие других. $\Rightarrow$ Для проверки совместности СЛАУ важно понятие ранга:
\begin{opr}
\textbf{Определение.} \emph{Ранг матрицы}~--- это
1)~количество ненулевых строк соответствующей матрицы ступенчатого вида;
2)~количество линейно независимых строк. \ 1)\,$\iff$\,2).
\end{opr}

\subsection{Вычисление обратной матрицы методом Гаусса}\label{ssec:obrgauss}
Пусть для $n\times n$-матрицы $A$ составлена <<расширенная>> матрица $(A\,|\,E)$ размерности
$n\times 2n$ (справа от $A$ добавлены столбцы единичной матрицы $E$). Выполняются элементарные
преобразования строк:
1)~умножение строки на число; 2)~перестановка двух строк; 3)~прибавление к строке другой строки,
умноженной на число. После этого получена матрица $(E\,|\,D)$ (т.\,е. элементарными
преобразованиями $A$ приведена к $E$, а $E$~--- к виду $D$). Тогда $D=A^{-1}$.

Для обоснования заметим, что указанные элементарные преобразования равносильны умножению $A$
\emph{слева} на некоторую матрицу $P$.
\begin{enumerate}[label=\arabic*)]
\item Пусть $P$~--- $n\times n$-матрица, у которой на диагонали $p_{ii}=\lambda$, остальные
диагональные элементы равны 1, элементы вне диагонали равны 0. Тогда по определению умножения
$PA$~--- это матрица $A$, в которой строка \textnumero\,$i$ умножена на $\lambda$, а остальные строки
не изменились.
\[
\prim
\begin{pmatrix}\lambda&0&0\\0&1&0\\0&0&1\end{pmatrix}
\begin{pmatrix}1&2&3\\0&-1&1\\4&0&-5\end{pmatrix}
=\begin{pmatrix}\lambda&2\lambda&3\lambda\\0&-1&1\\4&0&-5\end{pmatrix}.
\]
\item Пусть в $P=(p_{ij})$: \ $p_{ks}=p_{sk}=1$; \ $p_{ii}=1$ при $i\ne k$, $i\ne s$, а $p_{kk}=p_{ss}=0$;
остальные $p_{ij}=0$. Тогда $PA$ отличается от $A$ тем, что строки \textnumero\,$s$ и \textnumero\,$k$
поменялись местами.
\[
\prim
\begin{pmatrix}0&0&1\\0&1&0\\1&0&0\end{pmatrix}
\begin{pmatrix}1&2&3\\4&5&6\\7&8&9\end{pmatrix}
=\begin{pmatrix}7&8&9\\4&5&6\\1&2&3\end{pmatrix}.
\]
Здесь $p_{13}=p_{31}=1$, $p_{22}=1$, $p_{11}=p_{33}=0$~--- матрица перестановки строк 1 и 3 в
матрице $3\times3$.
\item Пусть в $P$, как и в $E$, на диагонали единицы, вне диагонали нули, кроме одного элемента
$p_{ks}=\lambda$. Тогда
\[
PA=\begin{pmatrix}
a_{11}&\dots&a_{1n}\\
\vdots&&\vdots\\
a_{k-1,1}&\dots&a_{k-1,n}\\
a_{k1}+\lambda a_{s1}&\dots&a_{kn}+\lambda a_{sn}\\
a_{k+1,1}&\dots&a_{k+1,n}\\
\vdots&&\vdots\\
a_{n1}&\dots&a_{nn}
\end{pmatrix}
\]
~--- на месте строки \textnumero\,$k$ теперь результат её сложения со строкой \textnumero\,$s$,
умноженной на $\lambda$.
\end{enumerate}

% =====================================================================
\section{Структура решения СЛАУ}
% =====================================================================
\[
Ax=b,\qquad A=(a_{ij})\ \ (n\times m),\quad b\in\mathbb R^n,\quad x\in\mathbb R^m.\tag{1}
\]
Соответствующая \emph{однородная} система:
\[
Ax=O_n.\tag{2}
\]
\q Какова структура решения СЛАУ (1)?

\begin{teor}{1}
Если $x_0$~--- решение (2), а $\tilde x$~--- решение (1), то $x_0+\tilde x$~--- решение (1).
\end{teor}

\zam В данном случае и вообще решение любого линейного уравнения или системы (алгебраического,
дифференциального и т.\,д.~--- уравнения, где над переменной выполняются лишь линейные операции)~---
это сумма: решение соответствующего однородного уравнения + решение самого неоднородного уравнения.

По теореме 1, зная какое-либо $\tilde x$ и добавляя к нему различные $x_0$ (если решение не
единственное), в итоге получим все решения (1).

\textbf{Доказательство.} $A\tilde x=b$, \ $Ax_0=O_n$. Сложим: $Ax_0+A\tilde x=b+O_n$; по
дистрибутивности $A(x_0+\tilde x)=b$ $\Rightarrow$ $x_0+\tilde x$~--- решение (1), т.\,к. получилось
верное равенство.\qedr

\q Каков вид решения (2)? Возможные варианты:
\begin{enumerate}[label=\arabic*)]
\item $m=n$ и $\rang A=n$ $\Rightarrow$ $\det A\ne0$ $\Rightarrow$ существует единственное решение (1):
$\tilde x=A^{-1}b$ (см. линейное матричное уравнение); для (2): $x_0=A^{-1}O_n=O_n$.

Однородная СЛАУ (2) с невырожденной матрицей имеет лишь нулевое решение $\Rightarrow$ СЛАУ (1) имеет
единственное решение.
\item Если (1) совместна, то по критерию Кронекера--Капелли (КК)\footnote{\textit{Критерий
Кронекера--Капелли:} система $Ax=b$ совместна $\iff$ $\rang A=\rang\bar A$. (Формулировка добавлена
для удобства.)} есть ещё один вариант:
\[
\rang\bar A=\rang A=r<m\quad\text{(ранг меньше числа переменных)}
\]
$\Rightarrow$ матрица $A$ имеет $r$ линейно независимых строк $\Rightarrow$ при элементарных
преобразованиях только $r$ строк останутся ненулевыми, а остальные $(n-r)$ станут нулевыми.

\textbf{Вывод:} в случае 2) при обратном ходе метода Гаусса получаем систему из $r$ уравнений,
а переменных $m>r$ $\Rightarrow$ часть переменных (их $m-r$)\ispr{В рукописи~--- $(n-r)$. Свободных
переменных столько, сколько неизвестных сверх ранга: $m-r$.} останется <<свободными>>!
$\Rightarrow$ (2) имеет бесконечно много решений.
\q Каков вид общего решения (2) в этом случае?
\end{enumerate}

\begin{teor}{2}
Если $x^{(1)},\dots,x^{(s)}$~--- решения (2), то любая их линейная комбинация
$\sum\limits_{i=1}^{s}\alpha_ix^{(i)}$~--- решение (2).
\end{teor}

\textbf{Доказательство.} По условию $Ax^{(i)}=O_n$, $i=1,\dots,s$. Умножим каждое равенство на число
$\alpha_i$: $\alpha_iAx^{(i)}=\alpha_iO_n=O_n$; сложим и применим ассоциативность:
$\sum_iA(\alpha_ix^{(i)})=O_n$; по дистрибутивности $A\bigl(\sum_i\alpha_ix^{(i)}\bigr)=O_n$
$\Rightarrow$ $\sum_i\alpha_ix^{(i)}$~--- решение (2).\qedr

\textbf{Вывод:} решение однородной СЛАУ (2) естественно представить в виде линейной комбинации
некоторых решений (\emph{базисных решений}). Очевидно, эти решения должны быть линейно
независимыми, иначе их линейная комбинация может дать очевидное решение $O_n$. Итак, целесообразно
следующее определение.

\begin{opr}
\textbf{Определение.} \emph{Фундаментальная система решений} (ФСР) для (2) при $\rang A=r<m$~---
это любой набор из $k=m-r$ линейно независимых решений. Если эти решения записать по столбцам
в матрицу, получится \emph{фундаментальная матрица} (ФМ).
\end{opr}

\zam Если $x_0$~--- решение (2), то $\alpha x_0$~--- тоже решение (2) при любом $\alpha$, поэтому в
случае $r<m$ ФСР бесконечно много (из каждого решения можно получить бесконечно много других).
Однако если две однородные СЛАУ имеют одну и ту же ФСР, то эти СЛАУ совпадают (равносильны).

\q Всякий ли набор линейно независимых решений может восстановить однородную СЛАУ?

\begin{teor}{3}
Если $X$~--- ФМ системы (2) (размер $m\times k$, $k=m-r$) и $C=(c_1,\dots,c_k)^T$~--- вектор
произвольных постоянных, то любое решение (2) имеет вид
\[
x_0=X\cdot C.\tag{3}
\]
Это вектор, каждая компонента которого~--- линейная комбинация с произвольными коэффициентами
$c_i$ ($x_0$~--- линейная комбинация столбцов $X$).
\end{teor}

\textbf{Доказательство.} ФМ составлена из ФСР $\Rightarrow$ $AX=O_{n\times k}$, т.\,е. ФМ
удовлетворяет (2). Умножим справа на $C$: $A(XC)=O_{n\times k}C=O_n$ $\Rightarrow$ $XC$~--- решение (2)
при любом $C$. С другой стороны, если взять любое решение (2) и рассмотреть (3), то отсюда
однозначно найдём $C$. Значит, (3)~--- вид общего решения (2).\qedr

\textbf{Вывод:} выбирая различные $C$ в (3), получим любое решение (2) $\Rightarrow$ по теореме 1 верна
\begin{teor}{4}
Любое решение (1) (общее решение (1)) имеет вид
\[
x=XC+\tilde x\tag{4}
\]
($X$~--- нулевая, если $\rang A=m$, т.\,е. $k=0$).
\end{teor}

\prim Найдём общее решение (семейство всех решений) системы\ispr{В рукописи в первом уравнении
написано $-3x_3$, но в матрице и дальше везде используется $+3$. Только при $+3x_3$ первое уравнение
равно сумме второго и третьего, ранг равен 2 и получается найденное решение.}
\[
\begin{cases}
2x_1-x_2+3x_3+4x_4-x_5=0,\\
x_1+5x_2-x_3-x_4-x_5=0,\\
x_1-6x_2+4x_3+5x_4=0.
\end{cases}\tag{5}
\]

\begin{enumerate}[label=\arabic*.]
\item \textbf{Прямой ход метода Гаусса:}
\[
\bar A=\left(\begin{array}{rrrrr|r}2&-1&3&4&-1&0\\1&5&-1&-1&-1&0\\1&-6&4&5&0&0\end{array}\right)
\sim\left(\begin{array}{rrrrr|r}2&-1&3&4&-1&0\\1&-6&4&5&0&0\end{array}\right)
\]
(строка 2 равна разности строк 1 и 3, её вычёркиваем).
$\rang\bar A=\rang A=2\ (=r)<5\ (=m)$ $\Rightarrow$ система совместна, решений бесконечно много.
\item По теореме 3 ФСР содержит $5-2=3$ решения, ФМ будет матрицей $5\times3$.\ispr{В рукописи~---
<<$2\times3$>>. В ФМ столько строк, сколько неизвестных ($m=5$), и столько столбцов, сколько решений
в ФСР ($k=3$).} Количество <<рабочих>> уравнений равно рангу, остальные~--- <<лишние>>.
\item \textbf{Обратный ход метода Гаусса:} $x_1=c_1$, $x_2=c_2$, $x_3=c_3$ (свободные), тогда
\[
x_4=\tfrac15(-c_1+6c_2-4c_3),\qquad x_5=\tfrac15(6c_1+19c_2-c_3).
\]
Скалярная форма $\iff$ векторная форма:
\[
x=\begin{pmatrix}x_1\\x_2\\x_3\\x_4\\x_5\end{pmatrix}
=\underbrace{\begin{pmatrix}1&0&0\\0&1&0\\0&0&1\\-\frac15&\frac65&-\frac45\\[2pt]\frac65&\frac{19}5&-\frac15\end{pmatrix}}_{X\ =\ \text{ФМ}}
\begin{pmatrix}c_1\\c_2\\c_3\end{pmatrix}.
\]
\end{enumerate}

% =====================================================================
\section{Собственные значения и собственные векторы матрицы}
% =====================================================================
Пусть $A=(a_{ij})$~--- $n\times n$-матрица. При умножении $A$ на вектор получается новый вектор
такого же размера, но, возможно, с другими компонентами $\Rightarrow$ возможен особый случай:
получится вектор, компоненты которого пропорциональны данному, т.\,е. выполнено
\[
Ax=\lambda x.\tag{1}
\]
\begin{opr}
\textbf{Определение.} Если выполнено (1), то 1)~$\lambda$~--- \emph{собственное значение} $A$ (число);
2)~$x\ne O_n$~--- соответствующий $\lambda$ \emph{собственный вектор} $A$.
\end{opr}

\textbf{Задача:} найти для $A$ \emph{спектр} (набор всех собственных значений), а также соответствующие
собственные векторы.

(1) $\iff$ $Ax-\lambda x=O_n$; \ $\lambda x=\lambda Ex$; по дистрибутивности
\[
(A-\lambda E)\,x=O_n.\tag{1$'$}
\]
(Матрица умножается на столбец $x$ справа.\ispr{В рукописи записано $x(A-\lambda E)$; при умножении
матрицы на столбец множитель $x$ должен стоять справа.})

Однородная СЛАУ (1$'$) имеет решение $x\ne O_n$ (см. определение, п.\,2) $\iff$
\[
\det(A-\lambda E)=0.\tag{2}
\]
$\det(A-\lambda E)$~--- \emph{характеристический многочлен} $A$, $P_n(\lambda)$;
(2)~--- \emph{характеристическое уравнение} (ХУ).

\zam Почему многочлен?
\[
\begin{vmatrix}a_{11}-\lambda&a_{12}&\dots&a_{1n}\\a_{21}&a_{22}-\lambda&\dots&a_{2n}\\
\vdots&&\ddots&\vdots\\a_{n1}&a_{n2}&\dots&a_{nn}-\lambda\end{vmatrix}
\]
~--- это сумма (со знаками) произведений элементов, взятых по одному из каждой строки (столбца), т.\,е.
сумма степенных выражений от $\lambda$~--- многочлен степени $n$.

\textbf{Итак, чтобы решить задачу:} 1) находим все корни (2); 2) подставляем каждый корень в (1$'$)
и находим соответствующие $x$.

\prim $A=\begin{pmatrix}-1&-6\\2&6\end{pmatrix}$.
\begin{enumerate}[label=\arabic*.]
\item ХУ: $\begin{vmatrix}-1-\lambda&-6\\2&6-\lambda\end{vmatrix}=0\iff(-1-\lambda)(6-\lambda)+12=0
\iff\lambda^2-5\lambda+6=0$; \ $\lambda_1=2$, $\lambda_2=3$~--- спектр.
\item $\lambda_1=2$: \ $(A-\lambda_1E)x=\begin{pmatrix}-3&-6\\2&4\end{pmatrix}\begin{pmatrix}x_1\\x_2\end{pmatrix}
=\begin{pmatrix}0\\0\end{pmatrix}$ (строку 1 умножим на $\frac23$ и прибавим к строке 2~--- она
обнулится) $\Rightarrow$ $-3x_1-6x_2=0$ $\Rightarrow$ $x_1=-2x_2$.

Так как ищем какой-либо собственный вектор (очевидно, их бесконечно много), то из семейства
$x^{(1)}=\begin{pmatrix}-2x_2\\x_2\end{pmatrix}$ выберем, например, $x^{(1)}=\begin{pmatrix}-2\\1\end{pmatrix}$.

Аналогично для $\lambda_2=3$: \ $x^{(2)}=\begin{pmatrix}3\\-2\end{pmatrix}$.
\end{enumerate}

\zam В данном случае и вообще собственные векторы, соответствующие различным собственным
значениям, ЛНЗ. Например, здесь $x^{(1)},x^{(2)}$ ЛНЗ, т.\,к.
\[
\alpha_1x^{(1)}+\alpha_2x^{(2)}=O_2\iff
\begin{pmatrix}-2&3\\1&-2\end{pmatrix}\begin{pmatrix}\alpha_1\\\alpha_2\end{pmatrix}=\begin{pmatrix}0\\0\end{pmatrix}
\overset{\det=1\ne0}{\iff}\alpha_1=\alpha_2=0.
\]

\prim $A=\begin{pmatrix}2&-1\\1&2\end{pmatrix}$. \ ХУ:
$\begin{vmatrix}2-\lambda&-1\\1&2-\lambda\end{vmatrix}=0\Rightarrow(2-\lambda)^2+1=0
\iff(\lambda-2)^2=-1$\ \ \textbf{(?!)}

В обычном (действительном) поле смысла нет $\Rightarrow$ чтобы решить ХУ, нужно новое определение.

\zam В математике всегда так: если задача не решается по одному определению, то придумывается новое.

Обобщим правило $x^2=a\iff x=\pm\sqrt a$.
\begin{opr}
\textbf{Определение.} $i$~--- \emph{мнимая единица}~--- символ, определяемый равенством $i^2=-1$.
\end{opr}
Тогда $\lambda_{1,2}=2\pm i$ $\Rightarrow$ как искать собственные векторы? (Ответ~--- в \S\,\ref{sec:primery}.)

% =====================================================================
\section{Комплексные числа}\label{sec:kch}
% =====================================================================
\begin{hwbox}
Рукописные листы <<Комплексные числа>> (л.\,1--3), вставлены сюда по смыслу.
\end{hwbox}
\hwlist{лист 1}\ Мы пришли к выводу, что для решения квадратного уравнения с $D<0$ целесообразно ввести
мнимую единицу~$i$: \ $i^2=-1$ (обычно пишут $i=+\sqrt{-1}$; формально подходит и $i=-\sqrt{-1}$).

При этом $ax^2+bx+c=0$ всегда имеет 2 корня (возможно, равные!):
$x_{1,2}=\dfrac{-b\pm\sqrt D}{2a}$ $\Rightarrow$ при $D<0$
\[
x_{1,2}=\frac{-b\pm\sqrt{(-1)(-D)}}{2a}
=\underbrace{-\frac{b}{2a}}_{\in\mathbb R}\pm\underbrace{\sqrt{-1}}_{=\,i}\cdot\underbrace{\frac{\sqrt{-D}}{2a}}_{\in\mathbb R}
\]
$\Rightarrow$ для таких случаев введём
\begin{opr}
\textbf{Определение.} $z=\alpha+i\beta$~--- \emph{комплексное число} (КЧ), $\alpha,\beta\in\mathbb R$:
$\alpha=\Rez z$~--- вещественная (действительная) часть, $\beta=\Imz z$~--- мнимая часть.
\end{opr}
В частности, необходимость ввести КЧ возникла при решении задачи о собственных значениях матрицы.

\q Как оперировать с КЧ?
\begin{enumerate}[label=\arabic*.]
\item \textbf{Произведение КЧ на вещественное число:} $k(\alpha+i\beta)=(k\alpha)+i(k\beta)$, $k\in\mathbb R$~---
и $\Rez z$, и $\Imz z$ умножились на $k$.
\item \textbf{Сумма двух КЧ:} $(\alpha+i\beta)+(a+ib)=(\alpha+a)+i(\beta+b)$~--- сложились вещественные
части и сложились мнимые части.
\item \textbf{Произведение двух КЧ} (применим дистрибутивность):
\[
(\alpha+i\beta)(a+ib)=\alpha a+i\alpha b+i\beta a+i^2\beta b=(\alpha a-\beta b)+i(\alpha b+\beta a).
\]
\item \hwlist{лист 2}\ \textbf{Частное КЧ} $\dfrac{z_1}{z_2}=?$~--- сведём к применению свойства 3
$\Rightarrow$ в знаменателе желательно получить вещественное число: $z\cdot(?)\in\mathbb R$.
\[
(\alpha+i\beta)(a+ib)\overset{(3)}{=}(\alpha a-\beta b)+i(\beta a+\alpha b),\qquad
\beta a+\alpha b=0\ \text{~--- должно быть}
\]
$\Rightarrow$ выберем $a=\alpha$, $b=-\beta$ $\Rightarrow$ $a+ib=\alpha-i\beta$ $\Rightarrow$ введём
\begin{opr}
\textbf{Определение.} Для $z=\alpha+i\beta$ число $\bar z=\alpha-i\beta$ называется \emph{сопряжённым} КЧ.
\end{opr}
Итак, для вычисления частного КЧ умножим и разделим его на число, сопряжённое знаменателю:
\[
\frac{z_1}{z_2}=\frac{z_1\,\bar z_2}{z_2\,\bar z_2}.
\]
\prim
\[
\frac{2-i}{3-2i}=\frac{(2-i)(3+2i)}{(3-2i)(3+2i)}
=\frac{2\cdot3-i^2\cdot2-3i+4i}{3^2-(2i)^2}=\frac{8+i}{9+4}=\frac{8}{13}+\frac{1}{13}\,i
\]
(в знаменателе применили $(a-b)(a+b)=a^2-b^2$).

\imp В знаменателе получилась сумма квадратов $\Rez(3-2i)$ и $\Imz(3-2i)$.
\begin{opr}
\textbf{Определение.} Величину $|z|=\sqrt{\alpha^2+\beta^2}$ назовём \emph{модулем} КЧ $z$.
\end{opr}
\zam $|z|=\sqrt{z\,\bar z}=\sqrt{(\alpha+i\beta)(\alpha-i\beta)}=\sqrt{\alpha^2+\beta^2}$.
\end{enumerate}

\subsection{Корни многочлена}
\hwlist{лист 3}\ При вычислении спектра матрицы $A$, т.\,е. её собственных значений~--- корней ХУ
\[
P_n(\lambda)=\det(A-\lambda E)=0,\tag{1}
\]
заметим, что многочлен в (1) в общем случае раскладывается на множители вида
\begin{itemize}
\item $(\lambda-a)^k$~--- соответствует вещественному корню $\lambda=a$ кратности $k$, т.\,е. в спектре
он учитывается $k$ раз;
\item $(\lambda^2+p\lambda+q)^s$, \ $p^2-4q<0$~--- соответствует паре комплексно-сопряжённых корней
\[
\lambda=-\frac p2\pm i\sqrt{-\frac D4}\qquad
\Bigl(\alpha=\Rez\lambda=-\frac p2,\ \ \beta=\Imz\lambda=\sqrt{-\tfrac D4}\Bigr),
\]
каждый кратности $s$; такой множитель даёт $2s$ чисел спектра.
\end{itemize}

\textbf{Вывод.} Многочлен $P_n(\lambda)$ степени $n$ (и вообще любой многочлен степени $n$!) имеет
ровно $n$ корней с учётом кратностей и с учётом комплексных корней.

\prim $\lambda^3+8=0\iff(\lambda+2)(\lambda^2-2\lambda+4)=0$ $\Rightarrow$ $\lambda_1=-2$,
\ $\lambda_{2,3}=\dfrac{2\pm\sqrt{-12}}{2}=1\pm i\sqrt3$. Всего 3 корня~--- столько, какова степень
многочлена.

% =====================================================================
\section{Примеры вычисления собственных значений и собственных векторов матрицы}\label{sec:primery}
% =====================================================================
\begin{hwbox}
Рукописные листы <<Примеры вычисления собственных значений и собственных векторов матрицы>>
(л.\,1--4), вставлены сюда по смыслу: пример \cn{1} продолжает пример из \S\,6.
\end{hwbox}
\hwlist{лист 1}\ \cn{1}\ \ $A=\begin{pmatrix}2&-1\\1&2\end{pmatrix}$.
Найдём спектр $A$ и собственные векторы $A$.
\begin{enumerate}[label=\arabic*.]
\item Ранее из ХУ получили $\lambda_{1,2}=2\pm i$~--- собственные значения.

\zam $\lambda_1\lambda_2=(2-i)(2+i)=2^2+1^2=5=\det A$ (здесь $\lambda_2=\bar\lambda_1$).
В данном случае и вообще: \emph{определитель матрицы равен произведению всех её собственных
значений} (кратные учитываются в произведении несколько раз!).
\item Найдём собственные векторы для $\lambda_{1,2}$.

$\lambda_1=2-i$ $\Rightarrow$ составим СЛАУ $(A-\lambda_1E)x=O_2$:
\[
\begin{pmatrix}i&-1\\1&i\end{pmatrix}\begin{pmatrix}x_1\\x_2\end{pmatrix}=\begin{pmatrix}0\\0\end{pmatrix}
\gauss{\text{II}+i\cdot\text{I}}
\begin{pmatrix}i&-1\\0&0\end{pmatrix}\begin{pmatrix}x_1\\x_2\end{pmatrix}=\begin{pmatrix}0\\0\end{pmatrix}
\iff ix_1-x_2=0\iff x_2=ix_1.
\]
Так как $x$~--- собственный вектор $\iff$ $\alpha x$~--- собственный вектор ($\alpha$~--- число, может
быть и комплексным), то для определённости достаточно выбрать (и в этом примере, и вообще!)
какой-либо вектор для данного $\lambda$. Пусть $x_1=1$ $\Rightarrow$
$x^{(1)}=\begin{pmatrix}1\\i\end{pmatrix}$.

$\lambda_2=2+i$ $\Rightarrow$ СЛАУ для соответствующего собственного вектора:
$\begin{pmatrix}-i&-1\\1&-i\end{pmatrix}\begin{pmatrix}x_1\\x_2\end{pmatrix}=O_2\iff-ix_1-x_2=0$
$\Rightarrow$ $x^{(2)}=\begin{pmatrix}1\\-i\end{pmatrix}$.
\end{enumerate}

\hwlist{лист 2}\ \cn{2}\ \ $A=\begin{pmatrix}0&2&0\\2&0&0\\0&0&3\end{pmatrix}$.

\zam В данном случае и вообще: если матрица \emph{симметрическая} ($A^T=A$), то 1)~её собственные
значения вещественны; 2)~соответствующие собственные векторы попарно ортогональны, т.\,е. их
скалярные произведения равны нулю.
\begin{enumerate}[label=\arabic*.]
\item ХУ: $\det(A-\lambda E)=0\iff\begin{vmatrix}-\lambda&2&0\\2&-\lambda&0\\0&0&3-\lambda\end{vmatrix}=0$
\ [разложим по строке 3] $\iff(3-\lambda)(\lambda^2-4)=0$ $\iff$ собственные значения
$\lambda_1=3$, $\lambda_2=-2$, $\lambda_3=2$ (вещественные, т.\,к. $A^T=A$).
\item Найдём соответствующие собственные векторы $x^{(1)},x^{(2)},x^{(3)}$.

$\lambda_1=3$: \ $(A-\lambda_1E)x^{(1)}=O_3\iff$
\[
\begin{pmatrix}-3&2&0\\2&-3&0\\0&0&0\end{pmatrix}x^{(1)}=O_3
\gauss{\text{II}+\frac23\,\text{I}}
\begin{pmatrix}-3&2&0\\0&-\frac53&0\end{pmatrix}x^{(1)}=O_2
\gauss{\text{II}\cdot3}
\begin{cases}-3x_1+2x_2=0,\\-5x_2=0\end{cases}
\]
$\Rightarrow$ $x_1=x_2=0$, $x_3$~--- любое $\Rightarrow$ выберем $x^{(1)}=\begin{pmatrix}0\\0\\1\end{pmatrix}$.

$\lambda_2=-2$: \ СЛАУ для $x^{(2)}$:
\[
\begin{pmatrix}2&2&0\\2&2&0\\0&0&5\end{pmatrix}x^{(2)}=O_3
\gauss{\text{II}-\text{I}}
\begin{pmatrix}2&2&0\\0&0&0\\0&0&5\end{pmatrix}x^{(2)}=O_3
\gauss{\text{I}:2,\ \text{III}:5}
\begin{cases}x_1+x_2=0,\\x_3=0\end{cases}
\]
$\Rightarrow$ выберем $x^{(2)}=\begin{pmatrix}-1\\1\\0\end{pmatrix}$.
$\lambda_3=2$: \ $\begin{pmatrix}-2&2&0\\2&-2&0\\0&0&1\end{pmatrix}x^{(3)}=O_3$ $\Rightarrow$ выберем
$x^{(3)}=\begin{pmatrix}1\\1\\0\end{pmatrix}$.

При этом $\bigl(x^{(1)}\bigr)^Tx^{(2)}=\bigl(x^{(1)}\bigr)^Tx^{(3)}=\bigl(x^{(2)}\bigr)^Tx^{(3)}=0$.
\end{enumerate}

\hwlist{лист 3}\ \textbf{Замечание (!!!).} В данном случае и вообще: если $A$~--- симметрическая
матрица, то $U^TAU$~--- диагональная матрица, и на её диагонали стоят собственные значения $A$
($U^TAU$~--- \emph{ортогональное преобразование} $A$). Здесь $U$~--- матрица, столбцы которой~---
нормированные собственные векторы $A$; при этом $U^T=U^{-1}$, $\det U=\pm1$, т.\,е. $U$~---
\emph{ортогональная} матрица.\ispr{В рукописи~--- $\det U=1$. У ортогональной матрицы $\det U=\pm1$;
в этом примере $\det U=-1$.}

\emph{Нормировка} вектора $x=(x_1,\dots,x_n)^T$~--- это результат его умножения на обратную длину:
\[
e=\frac{1}{|x|}\,x,\qquad |x|=\sqrt{\textstyle\sum_{i=1}^{n}x_i^2}.
\]
\vspace{-2pt}(Под корнем~--- сумма квадратов координат.\ispr{В рукописи под корнем пропущен квадрат:
$\sqrt{\sum x_i}$.})

В данном примере нормируем $x^{(1)},x^{(2)},x^{(3)}$ $\Rightarrow$ получим
\[
e^{(1)}=\begin{pmatrix}0\\0\\1\end{pmatrix},\quad
e^{(2)}=\begin{pmatrix}-\frac{\sqrt2}{2}\\[2pt]\frac{\sqrt2}{2}\\[2pt]0\end{pmatrix},\quad
e^{(3)}=\begin{pmatrix}\frac{\sqrt2}{2}\\[2pt]\frac{\sqrt2}{2}\\[2pt]0\end{pmatrix}
\ \Rightarrow\
U=\begin{pmatrix}0&-\frac{\sqrt2}{2}&\frac{\sqrt2}{2}\\[2pt]0&\frac{\sqrt2}{2}&\frac{\sqrt2}{2}\\[2pt]1&0&0\end{pmatrix}
\]
~--- ортогональная матрица $\Rightarrow$
\[
U^TAU=
\begin{pmatrix}0&0&1\\[2pt]-\frac{\sqrt2}{2}&\frac{\sqrt2}{2}&0\\[2pt]\frac{\sqrt2}{2}&\frac{\sqrt2}{2}&0\end{pmatrix}
\begin{pmatrix}0&2&0\\2&0&0\\0&0&3\end{pmatrix}
\begin{pmatrix}0&-\frac{\sqrt2}{2}&\frac{\sqrt2}{2}\\[2pt]0&\frac{\sqrt2}{2}&\frac{\sqrt2}{2}\\[2pt]1&0&0\end{pmatrix}
=\begin{pmatrix}3&0&0\\0&-2&0\\0&0&2\end{pmatrix}
\]
~--- на диагонали $\lambda_1,\lambda_2,\lambda_3$ (как и по теории!). Матрица полученного вида~---
\emph{жорданова нормальная форма} (ЖНФ) матрицы $A$.

\hwlist{лист 4}\ \cn{3}\ \ $A=\begin{pmatrix}5&1\\0&5\end{pmatrix}$.

\zam Для треугольной матрицы (и для диагональной в том числе!) диагональные элементы~--- это её
собственные значения $\Rightarrow$ приведение матрицы к такому виду может быть полезно для вычисления
собственных значений, но пока нам не ясно, подойдут ли для этого прежние <<элементарные
преобразования>>.
\begin{enumerate}[label=\arabic*.]
\item ХУ: $\begin{vmatrix}5-\lambda&1\\0&5-\lambda\end{vmatrix}=0\iff(\lambda-5)^2=0$ $\Rightarrow$
$\lambda_{1,2}=5$ кратности 2.
\item СЛАУ для собственных векторов:
$\begin{pmatrix}0&1\\0&0\end{pmatrix}\begin{pmatrix}x_1\\x_2\end{pmatrix}=\begin{pmatrix}0\\0\end{pmatrix}
\iff\begin{cases}x_2=0,\\x_1\ \text{~--- любое}\end{cases}$ $\Rightarrow$ выберем
$x^{(1)}=\begin{pmatrix}1\\0\end{pmatrix}$.
\end{enumerate}
\zam В данном случае и вообще: количество собственных векторов должно быть равно количеству
собственных значений, однако имеет смысл брать только ЛНЗ векторы; но, как видим, в этом примере
(при наличии кратных собственных значений) с этим проблема!

% =====================================================================
\section{Геометрический смысл комплексного числа}
% =====================================================================
\lead{Проблема:} с помощью алгебраической формы удобно умножать, складывать, делить комплексные
числа, но возникает вопрос: как извлечь корень $\sqrt[n]{z}=?$

Это требуется для вычисления полного набора корней многочлена, например:
\[
x^3+1=0\iff x=\sqrt[3]{-1}
\]
$\Rightarrow$ вещественный корень $x_1=-1$, но по основной теореме алгебры уравнение должно иметь
3 корня. Как найти остальные?

\lead{Идея.} В математике для решения любой проблемы применяется <<переформулирование>>, до тех
пор пока не получим вычислительный результат $\Rightarrow$ найдём другие формы записи КЧ.

Очевидно, существует взаимно однозначное соответствие $z=x+iy\iff$ точка $(x,y)\in\mathbb R^2$.
\begin{center}
\begin{tikzpicture}[scale=1.0,>=Stealth]
 \draw[->] (-0.6,0)--(4.2,0) node[below] {$x=\Rez z$};
 \draw[->] (0,-0.6)--(0,3.0) node[left] {$y=\Imz z$};
 \coordinate (P) at (3,2);
 \draw[dashed] (P)--(3,0) node[below] {$x$};
 \draw[dashed] (P)--(0,2) node[left] {$y$};
 \draw[->,thick,blue] (0,0)--(P) node[midway,above left] {$r$};
 \fill (P) circle (1.6pt) node[above right] {$(x,y)\iff z$};
 \draw[->] (0.8,0) arc (0:33.7:0.8);
 \node at (1.05,0.3) {$\varphi$};
 \draw[decorate,decoration={brace,mirror,amplitude=4pt}] (0,-0.35)--(3,-0.35)
   node[midway,below=4pt] {$r\cos\varphi$};
 \draw[decorate,decoration={brace,amplitude=4pt}] (-0.75,0)--(-0.75,2)
   node[midway,left=4pt] {$r\sin\varphi$};
\end{tikzpicture}
\end{center}
При изображении комплексных чисел на декартовой плоскости $\mathbb R^2$ будем называть её
\emph{комплексной плоскостью}~$\mathbb C$.

$r=|z|=\sqrt{x^2+y^2}$~--- \emph{модуль} $z$ (по теореме Пифагора)~--- это длина отрезка с концами
$(0,0)$ и $(x,y)$; этот отрезок~--- радиус-вектор точки $(x,y)$. Т.\,е. модуль КЧ $z$~--- длина его
радиус-вектора.

$\varphi$~--- \emph{аргумент} КЧ: \ $\rtg\varphi=\dfrac yx$, \ $x\ne0$, \ или $\varphi=\pm\dfrac\pi2+2\pi n$
(при $x=0$).

Очевидно, $\varphi$ имеет множество значений (с учётом периода $\rtg$) $\Rightarrow$ выделим
\emph{главное значение аргумента} КЧ: в любом промежутке длиной $2\pi$, для определённости
\[
\boxed{\ \varphi=\arg z\in(-\pi;\pi]\ }
\]
\[
\arg z=\begin{cases}
\rarctg\dfrac yx, & x>0\ \ \text{(I или IV четверть)},\\[6pt]
\dfrac\pi2, & x=0,\ y>0,\\[6pt]
\pi+\rarctg\dfrac yx, & x<0,\ y\geqslant0\ \ \text{(II четверть)},\\[6pt]
-\pi+\rarctg\dfrac yx, & x<0,\ y<0\ \ \text{(III четверть)},\\[6pt]
-\dfrac\pi2, & x=0,\ y<0.
\end{cases}
\]
\vspace{-2pt}(В третьей строке нестрогое неравенство $y\geqslant0$.\ispr{В рукописи~--- $y>0$. Тогда
отрицательные вещественные числа ($x<0$, $y=0$) не попадают ни в одну строку, а для них
$\arg z=\pi$.})

\prim $z=1-i$: \ $|z|=\sqrt{1^2+(-1)^2}=\sqrt2$; \ $\Rez z=1>0$, $\Imz z=-1<0$ (IV четверть)
$\Rightarrow$ $\arg z=\rarctg\frac{-1}{1}=-\frac\pi4$.

\begin{minipage}[c]{0.56\linewidth}
\prim $z=-3-2i$: \ $|z|=\sqrt{13}$; \ $\Rez z=-3<0$, $\Imz z=-2<0$ $\Rightarrow$ III четверть
\[
\Rightarrow\ \arg z=-\pi+\rarctg\frac{-2}{-3}=-\pi+\rarctg\frac23.
\]
\end{minipage}\hfill
\begin{minipage}[c]{0.4\linewidth}
\centering
\begin{tikzpicture}[scale=0.62,>=Stealth]
 \draw[->] (-3.8,0)--(1.5,0) node[above] {\small$\Rez z$};
 \draw[->] (0,-2.8)--(0,1.4) node[left] {\small$\Imz z$};
 \draw[dashed] (-3,-2)--(-3,0) node[above] {\small$-3$};
 \draw[dashed] (-3,-2)--(0,-2) node[right] {\small$-2$};
 \draw[->,thick,blue] (0,0)--(-3,-2) node[midway,above left=-2pt] {\small$\sqrt{13}$};
 \fill (-3,-2) circle (2pt);
 \draw[->,red] (0.7,0) arc (0:-146.3:0.7);
 \node[red] at (1.25,-0.85) {\small$\arg z$};
\end{tikzpicture}
\end{minipage}

\medskip
\textbf{Вывод:} КЧ можно записать в \emph{тригонометрической форме}:
\[
\boxed{\ z=r(\cos\varphi+i\sin\varphi)\ }\tag{1}
\]
Применим формулу Эйлера:
\[
\boxed{\ e^{i\varphi}=\cos\varphi+i\sin\varphi\ }\tag{2}
\]
\begin{hwbox}
\textbf{Интересный факт:} \ $\displaystyle\sum_{n=0}^{+\infty}q^n=\frac{1}{1-q}$, \ $|q|<1$.
\end{hwbox}
Из (1), (2) $\Rightarrow$ \emph{показательная форма} КЧ:
\[
\boxed{\ z=r\,e^{i\varphi}\ }\tag{3}
\]
\textbf{Вывод:} с помощью (3) удобно возводить КЧ в любую степень:
$z^n=r^ne^{in\varphi}\overset{(1)}{\Longrightarrow}$
\[
\boxed{\ z^n=r^n(\cos n\varphi+i\sin n\varphi)\ }\quad\text{~--- формула Муавра.}\tag{4}
\]
Из (3) следует формула для корня степени $n$: $\sqrt[n]z=\sqrt[n]r\,e^{i\varphi/n}$, но учтём, что
$\varphi$ имеет периодичность, чтобы получить $n$ значений корня. Это требуется для правильного
решения алгебраического уравнения:
\[
\boxed{\ \sqrt[n]z=\sqrt[n]r\left(\cos\frac{\arg z+2\pi k}{n}+i\sin\frac{\arg z+2\pi k}{n}\right),
\quad k=0,1,\dots,n-1\ }\tag{5}
\]
Очевидно, с геометрической точки зрения все значения $\sqrt[n]z$~--- вершины правильного
$n$-угольника, вписанного в окружность с центром $(0,0)$ и радиусом $R=\sqrt[n]r$.

\subsection{Примеры}
\begin{enumerate}[label=\arabic*.]
\item $\displaystyle\frac{z_1}{z_2}=\frac{r_1e^{i\varphi_1}}{r_2e^{i\varphi_2}}=\frac{r_1}{r_2}\,e^{i(\varphi_1-\varphi_2)}$~---
при делении КЧ модуль частного равен частному модулей, аргумент частного~--- разности
аргументов.\ispr{В рукописи про аргумент тоже написано <<частному модулей>>~--- описка.}
\item $z_1z_2\overset{(3)}{=}r_1r_2\,e^{i(\varphi_1+\varphi_2)}$~--- модуль произведения равен
произведению модулей, аргумент произведения~--- сумме аргументов.
\item $(1-i\sqrt3)^{90}=\circledast$. Можно применить бином Ньютона
$(a+b)^n=\sum\limits_{k=0}^{n}C_n^ka^kb^{n-k}$, $C_n^k=\dfrac{n!}{k!\,(n-k)!}$, но в данном случае
получим 91 слагаемое (неудобно).

Применим формулу Муавра $\Rightarrow$ вычислим модуль и аргумент основания:
$r=|1-i\sqrt3|=2$; \ $\arg(1-i\sqrt3)$ (IV четверть) $=\rarctg\dfrac{-\sqrt3}{1}=-\dfrac\pi3$.
\[
\circledast\overset{(4)}{=}2^{90}\Bigl(\cos\bigl(\underbrace{90\cdot(-\tfrac\pi3)}_{n\varphi=-30\pi}\bigr)
+i\sin\bigl(90\cdot(-\tfrac\pi3)\bigr)\Bigr)=2^{90},
\]
т.\,к. $\cos(-30\pi)=1$, $\sin(-30\pi)=0$.

\item $z^4+1=0$ \ (a) \ $\iff z=\sqrt[4]{-1}$. По основной теореме алгебры (a) как алгебраическое
уравнение степени 4 должно иметь 4 корня с учётом кратности корней и с учётом комплексных корней.
Поэтому $\sqrt[4]{-1}$ должен иметь 4 значения.

Применим формулу (5) $\Rightarrow$ найдём модуль и главное значение аргумента:
$|-1|=1$, \ $\arg(-1)=\pi$ \ ($-1=-1+0\cdot i$).
\[
\sqrt[4]{-1}\overset{(5)}{=}\sqrt[4]1\cdot\left(\cos\frac{-\pi+2\pi k}{4}+i\sin\frac{-\pi+2\pi k}{4}\right),
\quad k=0,1,2,3\ \ (k=0,\dots,n-1,\ n=4).
\]
(Вместо $\pi$ подставлено $-\pi$.\ispr{Это не ошибка по существу: $-\pi$~--- тоже аргумент числа $-1$,
набор корней от этого не меняется, меняется только их нумерация. По формуле (5) с главным значением
$\pi$ получилось бы $z_0=\frac{\sqrt2}2+i\frac{\sqrt2}2$ и т.\,д.})

\begin{minipage}[c]{0.5\linewidth}
$k=0$: \ $z_0=\dfrac{\sqrt2}{2}-i\dfrac{\sqrt2}{2}$;\\[3pt]
$k=1$: \ $z_1=\dfrac{\sqrt2}{2}+i\dfrac{\sqrt2}{2}$;\\[3pt]
$k=2$: \ $z_2=-\dfrac{\sqrt2}{2}+i\dfrac{\sqrt2}{2}$;\\[3pt]
$k=3$: \ $z_3=-\dfrac{\sqrt2}{2}-i\dfrac{\sqrt2}{2}$.

\medskip
Изобразим эти числа на комплексной плоскости $\mathbb C$: они~--- вершины квадрата.
\end{minipage}\hfill
\begin{minipage}[c]{0.44\linewidth}
\centering
\begin{tikzpicture}[scale=1.25,>=Stealth]
 \draw[->] (-1.5,0)--(1.6,0) node[below] {\small$\Rez z$};
 \draw[->] (0,-1.4)--(0,1.5) node[left] {\small$\Imz z$};
 \draw (0,0) circle (1);
 \node[below left] at (1,0) {\scriptsize$1$};
 \node[above left] at (0,1) {\scriptsize$1$};
 \node[below left] at (-1,0) {\scriptsize$-1$};
 \node[below left] at (0,-1) {\scriptsize$-1$};
 \draw[thick,blue] (45:1)--(135:1)--(225:1)--(315:1)--cycle;
 \foreach \a/\n/\p in {45/z_1/above right,135/z_2/above left,225/z_3/below left,315/z_0/below right}{
   \draw[dashed,gray] (0,0)--(\a:1); \fill (\a:1) circle (1.5pt) node[\p] {\small$\n$};}
 \draw[->] (0.35,0) arc (0:45:0.35); \node at (0.55,0.2) {\tiny$\frac\pi4$};
\end{tikzpicture}
\end{minipage}

\item \textbf{Решим квадратное уравнение} \ $z^2+(-3+2i)z+5-i=0$.
\[
D=(-3+2i)^2-4(5-i)=9-12i+4i^2-20+4i=-15-8i\ \Rightarrow\ \sqrt D=?
\]
Можно применить формулу (5) для $\sqrt[n]z$ и выбрать одно из значений, т.\,к. в формуле корней и так
стоит <<$\pm\sqrt D$>>. Но при этом $|D|=\sqrt{289}=17$, \ $\arg D=-\pi+\rarctg\frac8{15}$~--- не
подходит: по формуле для корня получим $\cos\bigl(-\frac\pi2+\frac12\rarctg\frac8{15}+\pi k\bigr)$
или $\sin(\ldots)$~--- неудобно (можно сделать проще).

\textbf{Вычислим корень иначе} $\Rightarrow$ идея: корень удобно извлекать из квадрата $\Rightarrow$
применим МНК (метод неопределённых коэффициентов):
\[
D=(\alpha+i\beta)^2\iff-15-8i=(\alpha^2-\beta^2)+i\,(2\alpha\beta)
\ \Rightarrow\
\begin{cases}\alpha^2-\beta^2=-15,\\2\alpha\beta=-8\end{cases}
\]
$\Rightarrow$ $\beta=-\dfrac4\alpha$ $\Rightarrow$ $\alpha^2-\dfrac{16}{\alpha^2}+15=0$.
Отсюда $\alpha^4+15\alpha^2-16=0$ $\Rightarrow$ $\alpha^2=1$ (корень $\alpha^2=-16$ не подходит, т.\,к.
$\alpha\in\mathbb R$) $\Rightarrow$ $\alpha=\pm1$.

Выберем одно из решений: $\alpha=1$, $\beta=-4$ $\Rightarrow$ $D=(1-4i)^2$ $\Rightarrow$ $\sqrt D=1-4i$.

Имеем корни\ispr{В рукописи из $\alpha^4+15\alpha^2-16=0$ получено $\alpha^2=16$ или $-1$ (перепутаны
знаки корней), откуда $\alpha=4$, $\beta=-1$, $\sqrt D=4-i$. Но $(4-i)^2=15-8i\ne D$, поэтому и корни
$\frac{7-3i}2$, $\frac{-1-i}2$ в рукописи неверные. Верные корни проверены подстановкой в уравнение.}:
\[
z_{1,2}=\frac{3-2i\pm(1-4i)}{2}:\qquad z_1=\frac{4-6i}{2}=2-3i,\qquad z_2=\frac{2+2i}{2}=1+i.
\]
Проверка (теорема Виета): $z_1+z_2=3-2i=-(-3+2i)$ \checkmark, \ $z_1z_2=(2-3i)(1+i)=5-i$ \checkmark.

\item На $\mathbb C$ множество определено неравенством $|z-a|>R$.

\imp Для множества основная характеристика~--- это \emph{граница}; для множества, определённого
неравенством, граница определяется соответствующим равенством.

\begin{minipage}[c]{0.55\linewidth}
В данном случае граница~--- $|z-a|=R$ (расстояние от точки $z$ до заданной точки $a$ равно $R$)
$\Rightarrow$ это уравнение окружности с центром $a$ и радиусом $R$. Граница разделила $\mathbb C$
на 2 части: внутри и вне окружности. Выберем нужную часть по аналогии с методом интервалов:
проверим точку $a$ (внутри): $|a-a|=0>R$~--- неверно $\Rightarrow$ внутренняя часть не подходит
$\Rightarrow$ подходит внешняя часть~--- это всё остальное (окружность пунктиром: неравенство
строгое).
\end{minipage}\hfill
\begin{minipage}[c]{0.41\linewidth}
\centering
\begin{tikzpicture}[scale=0.8,>=Stealth]
 \begin{scope}
  \clip (-1.2,-1.4) rectangle (4.2,3.2);
  \fill[pattern=north east lines,pattern color=blue!60,even odd rule]
    (-1.2,-1.4) rectangle (4.2,3.2) (1.8,1.3) circle (1.1);
 \end{scope}
 \draw[->] (-1.2,0)--(4.3,0) node[below] {\small$x$};
 \draw[->] (0,-1.4)--(0,3.3) node[left] {\small$y$};
 \draw[dashed,thick] (1.8,1.3) circle (1.1);
 \fill (1.8,1.3) circle (1.6pt) node[above] {\small$a$};
 \draw[<->] (1.8,1.3)--++(-140:1.1) node[midway,above left=-2pt] {\small$R$};
\end{tikzpicture}
\end{minipage}
\end{enumerate}

% =====================================================================
\section{Аналитическая геометрия}
% =====================================================================
\lead{Идея:} Аполлоний (Др.~Греция). В новое время идею возродил П.~Ферма.
Основа: метод координат (Р.~Декарт, <<Геометрия>>).

\begin{opr}
\textbf{Определение.} \emph{Аналитическая геометрия} (АГ)~--- это раздел математики, в котором
геометрические свойства линий и поверхностей изучаются при помощи алгебраических уравнений на
основе метода координат.
\end{opr}

\subsection{Понятие вектора}
Понятие <<вектор>> (от лат. <<переносящий>>) предложил Гамильтон, 1853\,г. Современное понятие
предложили Хевисайд и Гиббс.

\begin{opr}
\textbf{Определение.} Если на прямой выбраны две точки $A$, $B$ и указано, что $A$~--- начальная, $B$~---
конечная, то определён \emph{направленный отрезок} (НО) $\overline{AB}$.
\end{opr}
\begin{center}
\begin{tikzpicture}[>=Stealth]
 \draw (-0.3,-0.15)--(3.3,1.65) node[right] {$a$};
 \draw[->,very thick] (0.4,0.2)--(2.4,1.2);
 \fill (0.4,0.2) circle (1.6pt) node[below right] {$A$};
 \node[above left] at (2.4,1.2) {$B$};
 \node at (4.6,0.4) {$\overline{AB}$~--- обозначение};
\end{tikzpicture}
\end{center}

\zam Зачем НО? В физике, в технике многие величины (например, сила $\vec F$) характеризуются не только
величиной, но и направлением $\Rightarrow$ их естественная характеристика~--- это НО.

\begin{minipage}[c]{0.6\linewidth}
Например, если материальная точка под действием силы $\vec F$ вращается вокруг оси, то возникает
понятие момента вращения (момента силы), и его характеризуют 2 вектора: плечо $\vec r$ и сила $\vec F$.
\end{minipage}\hfill
\begin{minipage}[c]{0.36\linewidth}
\centering
\begin{tikzpicture}[>=Stealth,scale=0.9]
 \draw[dashed] (0,-0.3)--(0,2.4) node[above] {\small ось};
 \draw[->,thick] (0,1.6)--(1.4,0.6) node[midway,above right=-2pt] {\small$\vec r$ (плечо)};
 \fill (1.4,0.6) circle (1.8pt) node[below] {\small точка};
 \draw[->,thick,red] (1.4,0.6)--(2.4,1.4) node[right] {\small$\vec F$};
\end{tikzpicture}
\end{minipage}

\begin{opr}
\textbf{Определение.} \emph{Длина} НО $\overline{AB}$~--- это длина соответствующего отрезка прямой $|AB|$.
\end{opr}
\begin{opr}
\textbf{Определение.} Два НО \emph{коллинеарны} ($\parallel$), если они лежат на параллельных прямых:
\ $\overline{AB}\parallel\overline{CD}\iff AB\parallel CD$.
\end{opr}

\zam По сути мы неявно используем аксиомы евклидовой геометрии, которые описывают свойства
основных, не определяемых понятий: 1)~точка; 2)~прямая; 3)~плоскость.

\begin{opr}
\textbf{Определение.} Будем говорить, что $\overline{AB}$ и $\overline{CD}$ \emph{одинаково
(противоположно) направлены}, и обозначать $\overline{AB}\upuparrows\overline{CD}$
($\overline{AB}\uparrow\downarrow\overline{CD}$), если $\overline{AB}\parallel\overline{CD}$ и концы $B$ и $D$
лежат в одной полуплоскости (в разных полуплоскостях) относительно прямой $AC$, если плоскость
опирается на прямые $AB$ и $CD$.
\end{opr}
\begin{center}
\begin{tikzpicture}[>=Stealth,scale=0.85]
 \begin{scope}
  \draw (-0.4,-0.3)--(2.6,1.95);
  \draw (0.8,-1.6)--(3.8,0.65);
  \coordinate (A) at (0,0); \coordinate (B) at (1.2,0.9);
  \coordinate (C) at (1.6,-1.0); \coordinate (D) at (2.8,-0.1);
  \draw[->,very thick] (A)--(B); \draw[->,very thick] (C)--(D);
  \draw[dashed] ($(A)!-0.4!(C)$)--($(A)!1.6!(C)$) node[below] {\small$AC$};
  \foreach \P in {A,B,C,D}{\fill (\P) circle (1.5pt);}
  \node[left] at (A) {$A$}; \node[above left] at (B) {$B$};
  \node[below left] at (C) {$C$}; \node[above left] at (D) {$D$};
  \node at (1.6,2.4) {$\overline{AB}\upuparrows\overline{CD}$};
 \end{scope}
 \begin{scope}[xshift=7cm]
  \draw (-0.4,-0.3)--(2.6,1.95);
  \draw (0.8,-1.6)--(3.8,0.65);
  \coordinate (A) at (0,0); \coordinate (B) at (1.2,0.9);
  \coordinate (C) at (2.6,-0.25); \coordinate (D) at (1.4,-1.15);
  \draw[->,very thick] (A)--(B); \draw[->,very thick] (C)--(D);
  \draw[dashed] ($(A)!-0.3!(C)$)--($(A)!1.4!(C)$) node[right] {\small$AC$};
  \foreach \P in {A,B,C,D}{\fill (\P) circle (1.5pt);}
  \node[left] at (A) {$A$}; \node[above left] at (B) {$B$};
  \node[above right] at (C) {$C$}; \node[below left] at (D) {$D$};
  \node at (1.6,2.4) {$\overline{AB}\uparrow\downarrow\overline{CD}$};
 \end{scope}
\end{tikzpicture}
\end{center}

\begin{opr}
\textbf{Определение.} $\vec a=\vec b\iff\begin{cases}1)\ |\vec a|=|\vec b|,\\2)\ \vec a\upuparrows\vec b.\end{cases}$
\end{opr}
\begin{opr}
\textbf{\color{red!70!black}!}\ \textbf{Определение.} \emph{Вектор}~--- это множество всех равных
направленных отрезков.
\end{opr}

\zam Зачем вектор? Это понятие позволяет выбирать удобный НО из множества равных и прикладывать его
к любой начальной точке. Действительно, для любого $\overline{AB}$ и любой точки $A_1$ существует
единственный $\overline{A_1B_1}=\overline{AB}$, который строим так:

\begin{minipage}[c]{0.58\linewidth}
\begin{enumerate}[label=\cn{\arabic*}]
\item строим прямую $l\parallel AB$ так, что $A_1\in l$;
\item строим $AA_1$;
\item в полуплоскости, где лежит $B$, выбираем $B_1\in l$: $|AB|=|A_1B_1|$; при этом
$\overline{A_1B_1}\upuparrows\overline{AB}$.
\end{enumerate}
Итак, $\overline{A_1B_1}=\overline{AB}$ по определению.
\end{minipage}\hfill
\begin{minipage}[c]{0.38\linewidth}
\centering
\begin{tikzpicture}[>=Stealth,scale=0.85]
 \coordinate (A) at (0,0); \coordinate (B) at (1.4,1.9);
 \coordinate (A1) at (2.0,-0.3); \coordinate (B1) at (3.4,1.6);
 \draw ($(A)!-0.25!(B)$)--($(A)!1.15!(B)$);
 \draw[->,very thick] (A)--(B);
 \draw[dashed] ($(A1)!-0.5!(B1)$)--($(A1)!1.35!(B1)$) node[right] {\small$l$};
 \draw[dashed,gray] (A)--(A1);
 \draw[->,very thick,blue] (A1)--(B1);
 \fill (A) circle (1.5pt) node[left] {$A$}; \node[above left] at (B) {$B$};
 \fill (A1) circle (1.5pt) node[below right] {$A_1$}; \node[above left] at (B1) {$B_1$};
\end{tikzpicture}
\end{minipage}

\subsection{Линейные операции над векторами}
\begin{opr}
\textbf{Определение.} \emph{Сумма векторов}:
\[
\overrightarrow{AB}+\overrightarrow{BC}=\overrightarrow{AC}\tag{1}
\]
~--- первый вектор представлен НО, к его концу прикладываем начало второго вектора в виде НО.
\end{opr}

\zam (1) с геометрической точки зрения правильно называть \emph{правилом треугольника}. Если выбрать
$B_1$: $\overline{AB_1}\upuparrows\overline{BC}$ и $|AB_1|=|BC|$, то $\overline{AB_1}=\overline{BC}$
(по определению), тогда сложение векторов реализует \emph{правило параллелограмма}:
\begin{center}
\begin{tikzpicture}[>=Stealth,scale=0.95]
 \begin{scope}
  \coordinate (A) at (0,0); \coordinate (B) at (1.6,1.1); \coordinate (C) at (3,0.2);
  \draw[->,thick] (A)--(B) node[midway,above left] {\small$\overrightarrow{AB}$};
  \draw[->,thick] (B)--(C) node[midway,above right] {\small$\overrightarrow{BC}$};
  \draw[->,thick,blue] (A)--(C) node[midway,below] {\small$\overrightarrow{AC}$};
  \node[left] at (A) {$A$}; \node[above] at (B) {$B$}; \node[right] at (C) {$C$};
  \node at (1.5,-1.0) {правило треугольника};
 \end{scope}
 \begin{scope}[xshift=6.5cm]
  \coordinate (A) at (0,0); \coordinate (B) at (1.6,1.1); \coordinate (C) at (3,0.2);
  \coordinate (B1) at (1.4,-0.9);
  \draw[->,thick] (A)--(B); \draw[->,thick] (A)--(B1);
  \draw[dashed] (B)--(C); \draw[dashed] (B1)--(C);
  \draw[->,thick,blue] (A)--(C);
  \node[left] at (A) {$A$}; \node[above] at (B) {$B$}; \node[right] at (C) {$C$};
  \node[below] at (B1) {$B_1$};
  \node at (1.5,-1.6) {правило параллелограмма};
 \end{scope}
\end{tikzpicture}
\end{center}
Если складывается несколько векторов~--- получим \emph{правило многоугольника}:
\[
\overrightarrow{A_1A_2}+\overrightarrow{A_2A_3}+\dots+\overrightarrow{A_{n-1}A_n}=\overrightarrow{A_1A_n}.
\]

Сложение векторов обладает обычными алгебраическими свойствами:
\begin{enumerate}[label=\cn{\arabic*}]
\item \textbf{коммутативность}: $\vec a+\vec b=\vec b+\vec a$. Здесь учтена возможность оперировать
векторами через удобные НО, входящие в эти векторы.
\item \textbf{ассоциативность}: $\vec a+(\vec b+\vec c)=(\vec a+\vec b)+\vec c$. Иллюстрация~--- самостоятельно
(рисунок справа добавлен в это место, оставленное в конспекте пустым).
\item \textbf{наличие нейтрального элемента}: $\vec a+\vec0=\vec a$; его представляем направленным
отрезком с равными концами: $\vec0=\overline{AA}$, т.\,к. для любого $\overline{BA}$:
$\overline{BA}+\overline{AA}\overset{(1)}{=}\overline{BA}$.
\end{enumerate}
\begin{center}
\begin{tikzpicture}[>=Stealth,scale=0.95]
 \begin{scope}
  \coordinate (O) at (0,0); \coordinate (P) at (2.6,1.0); \coordinate (Q) at (0.9,1.6);
  \coordinate (S) at ($(P)+(Q)$);
  \draw[->,thick] (O)--(P) node[midway,below] {\small$\vec a$};
  \draw[->,thick] (P)--(S) node[midway,right] {\small$\vec b$};
  \draw[->,thick,gray] (O)--(Q) node[midway,left] {\small$\vec b$};
  \draw[->,thick,gray] (Q)--(S) node[midway,above] {\small$\vec a$};
  \draw[->,thick,blue] (O)--(S) node[pos=0.55,above left=-2pt] {\small$\vec c$};
  \node at (1.8,-0.9) {\small$\vec c=\vec a+\vec b=\vec b+\vec a$};
 \end{scope}
 \begin{scope}[xshift=7cm]
  \coordinate (O) at (0,0); \coordinate (A) at (1.8,-0.2); \coordinate (B) at (3.0,1.1);
  \coordinate (C) at (1.6,2.2);
  \draw[->,thick] (O)--(A) node[midway,below] {\small$\vec a$};
  \draw[->,thick] (A)--(B) node[midway,right] {\small$\vec b$};
  \draw[->,thick] (B)--(C) node[midway,above right] {\small$\vec c$};
  \draw[->,dashed,teal] (O)--(B) node[pos=0.72,above=1pt] {\scriptsize$\vec a+\vec b$};
  \draw[->,dashed,orange!80!black] (A)--(C) node[pos=0.25,left] {\scriptsize$\vec b+\vec c$};
  \draw[->,thick,blue] (O)--(C);
  \node[left] at (O) {$O$}; \node[right] at (A) {$A$}; \node[right] at (B) {$B$}; \node[above] at (C) {$C$};
  \node[align=center] at (1.6,-1.05) {\small$(\vec a+\vec b)+\vec c=\overrightarrow{OB}+\overrightarrow{BC}=\overrightarrow{OC}$\\
   \small$\vec a+(\vec b+\vec c)=\overrightarrow{OA}+\overrightarrow{AC}=\overrightarrow{OC}$};
 \end{scope}
\end{tikzpicture}
\end{center}

\begin{opr}
\textbf{Определение.} \emph{Умножение вектора на число.} Будем говорить, что $\vec b=\alpha\vec a$
($\alpha$~--- число), если
1)~$|\vec b|=|\alpha|\cdot|\vec a|$; \
2)~$\vec b\parallel\vec a$, причём $\vec b\upuparrows\vec a$ при $\alpha>0$ или $\vec b\uparrow\downarrow\vec a$
при $\alpha<0$.
\end{opr}
\begin{center}
\begin{tikzpicture}[>=Stealth,scale=0.9]
 \draw (-1.5,-1.2)--(1.6,1.28) node[right] {$l$};
 \fill (0,0) circle (1.6pt);
 \draw[->,very thick] (0,0)--(0.8,0.64) node[above left] {$\vec a$};
 \draw[->,very thick,red] (0,0)--(-1.3,-1.04) node[above left] {$\vec b$};
 \node at (2.7,0) {$\alpha<0$};
\end{tikzpicture}
\end{center}
\textbf{Свойства:} 1)~дистрибутивность: $(\alpha+\beta)\vec a=\alpha\vec a+\beta\vec a$, \
$\alpha(\vec a+\vec b)=\alpha\vec a+\alpha\vec b$~--- проиллюстрировать самостоятельно;
\ 2)~наличие нейтрального элемента: $1\cdot\vec a=\vec a$.

\q Какой элемент противоположный при сложении векторов; как определить разность векторов
$\vec a-\vec b=?$ (+ рисунок).

\lead{Идея.} Если с несколькими векторами совершить линейные операции, то получим их \emph{линейную
комбинацию} $\sum_{k=1}^{n}\alpha_k\vec a_k$ (ЛК). При этом ясно, что линейная комбинация
векторов~--- тоже вектор $\Rightarrow$ идея: найти такой набор векторов, через который в виде ЛК можно
представить любой вектор $\Rightarrow$ при этом информация о векторе станет набором числовых
коэффициентов.

\begin{opr}
\textbf{Определение.} Множество векторов с операциями умножения на число и сложения~--- это
\emph{векторное пространство} (линейное пространство).
\end{opr}

\subsection{Базис}
Чтобы реализовать предыдущую идею, учтём, что естественно выбирать максимальный набор ЛНЗ векторов,
чтобы через их ЛК представлять другие векторы $\Rightarrow$ введём
\begin{opr}
\textbf{Определение.} \emph{Базис} в векторном пространстве~--- это упорядоченный набор векторов:
\begin{enumerate}[label=\arabic*)]
\item на прямой~--- один произвольно выбранный ненулевой вектор;
\item на плоскости~--- любые 2 неколлинеарных вектора ($\vec a$ и $\vec b$: $\vec a\nparallel\vec b$);
\item в пространстве~--- любые 3 некомпланарных вектора (не существует представителей их НО,
лежащих в одной плоскости).
\end{enumerate}
\end{opr}
К пункту 2): по определению умножения вектора на число
\begin{multline*}
\forall\alpha\in\mathbb R:\ \vec a\ne\alpha\vec b\iff\vec a-\alpha\vec b\ne\vec0\iff{}\\
\iff\forall\alpha_1,\alpha_2\ (\text{не равных 0 одновременно}):\ \alpha_1\vec a+\alpha_2\vec b\ne\vec0
\iff\vec a,\vec b\ \text{~--- ЛНЗ}.
\end{multline*}
\zam По аналогии с пунктом 2) можно доказать, что некомпланарные векторы ЛНЗ.

Очевидно, базис можно выбрать неоднозначно. В частности, если базисные векторы 1)~имеют длину 1,
т.\,е. нормированы; 2)~попарно ортогональны,~--- то такой базис называют \emph{ортонормированным},
или \emph{декартовым}. Он состоит из векторов $\vec i,\vec j,\vec k$:
\[
|\vec i|=|\vec j|=|\vec k|=1,\qquad \vec i\perp\vec j,\quad \vec i\perp\vec k,\quad \vec j\perp\vec k.
\]
Все иные базисы~--- \emph{аффинные} (косоугольные).

\zam Так как в основе понятия <<вектор>> лежит НО, то целесообразно задать базис через точки.
\begin{opr}
\textbf{Определение.} \emph{Репер}:
1)~на прямой~--- две точки $A$, $B$ $\Rightarrow$ $\overrightarrow{AB}\ne\vec0$~--- базис;
2)~на плоскости~--- любые точки $O,A,B$: $\overrightarrow{OA}\nparallel\overrightarrow{OB}$ $\Rightarrow$
$\overrightarrow{OA},\overrightarrow{OB}$~--- базис;
3)~в пространстве~--- какие-то 4 точки $O,A,B,C$, не лежащие в одной плоскости $\Rightarrow$
$\overrightarrow{OA},\overrightarrow{OB},\overrightarrow{OC}$~--- базис.
\end{opr}
\begin{center}
\begin{tikzpicture}[>=Stealth,scale=0.9]
 \begin{scope}
  \draw (-0.3,-0.15)--(2.9,1.45) node[right] {$a$};
  \draw[->,very thick] (0.3,0.15)--(2.1,1.05);
  \fill (0.3,0.15) circle (1.5pt) node[below right] {$A$}; \node[above left] at (2.1,1.05) {$B$};
  \node at (1.3,-0.7) {\small 1) прямая};
 \end{scope}
 \begin{scope}[xshift=5cm]
  \draw[->,very thick] (0,0)--(2.2,0) node[right] {$A$};
  \draw[->,very thick] (0,0)--(1.2,1.3) node[above] {$B$};
  \fill (0,0) circle (1.5pt) node[left] {$O$};
  \draw (0.5,0) arc (0:47:0.5); \node at (0.85,0.35) {\scriptsize$\alpha\ne0$};
  \node at (1.1,-0.7) {\small 2) плоскость};
 \end{scope}
 \begin{scope}[xshift=10cm]
  \draw[->,very thick] (0,0)--(2.2,-0.8) node[right] {$A$};
  \draw[->,very thick] (0,0)--(2.3,0.3) node[right] {$B$};
  \draw[->,very thick] (0,0)--(1.0,1.4) node[above] {$C$};
  \fill (0,0) circle (1.5pt) node[left] {$O$};
  \node at (1.1,-1.2) {\small 3) пространство};
 \end{scope}
\end{tikzpicture}
\end{center}

\begin{opr}
\textbf{Определение.} Пусть $\vec e_1,\vec e_2,\vec e_3$~--- базис и для вектора $\vec a$:
$\vec a=\alpha_1\vec e_1+\alpha_2\vec e_2+\alpha_3\vec e_3$; тогда набор чисел
$(\alpha_1,\alpha_2,\alpha_3)$~--- это \emph{координаты} $\vec a$ в данном базисе.
\end{opr}
\begin{teor}{1}
Любой вектор в данном базисе имеет единственный набор координат.
\end{teor}

\end{document}
